\chapter{Linear Algebra, Revisited}

\begin{center}\itshape
Why is $F_{n+1}F_{n-1}-F_n^2$ always $\pm1$? Why is the determinant of $\left(\frac1{i+j-1}\right)_{i,j}$ positive, though tiny?\\
Matrices answer questions that do not mention matrices.
\end{center}

A nineteenth-century puzzle, found among the papers of Lewis Carroll, cuts an $8\times8$ chessboard into four pieces and reassembles them into a
$5\times13$ rectangle (Figure~\ref{fig:chessboard}). The chessboard has $64$ squares, the rectangle $65$. Where does the extra square come from? The
pieces do not quite fit: along the diagonal of the rectangle runs a long, thin parallelogram of area exactly $1$, too thin to see at the scale of a
drawing, since it is at most a fifth of a square wide. The numbers $5$, $8$, $13$ are consecutive Fibonacci numbers, and the area of the gap is
$5\cdot13-8^2=1$, an instance of the identity of the epigraph, which Jean-Dominique Cassini, the astronomer who measured the solar system, published in
1680, and which Robert Simson rediscovered in 1753. That identity is the determinant of a $2\times2$ matrix in disguise, and the gap is a determinant
read as an area: the two themes of this chapter.

Determinants were discovered twice, on opposite sides of the world and within ten years. In Japan, Seki Takakazu wrote in 1683 a manuscript on the
elimination of unknowns from systems of equations in which the $3\times3$ determinant appears, with a rule for its signs; in Hanover, Leibniz
described the same quantity, with a notation for the entries that is essentially our $a_{ij}$, in a letter to l'H\^opital of 1693, which was not
printed until 1850. Gabriel Cramer gave his rule for solving linear systems in 1750; Cauchy, in a memoir of 1812, gave the general definition with the
signs of permutations, proved that the determinant of a product is the product of the determinants, and fixed the name. The word \emph{matrix} is
Sylvester's, from 1850, and the theory of matrices as objects that one adds and multiplies is Cayley's, from a memoir of 1858 in which he also stated,
and checked for $2\times2$ and $3\times3$, the theorem that a matrix satisfies its own characteristic equation; we prove it in full below, as
Theorem~\ref{thm:cayleyhamilton}. That a real symmetric matrix has real eigenvalues was proved by Cauchy in 1829, not for matrices, which did not yet
exist, but for the axes of an ellipsoid; the Hilbert matrix of the epigraph is from a paper of 1894 on approximation by polynomials; and the most
valuable eigenvector in history, the one that ranks the pages of the web, was computed in 1998 by two students at Stanford. This chapter revisits
linear algebra with an eye for such disguises: determinants, eigenvalues and diagonalisation, used to answer questions that never mention a matrix.

\begin{figure}[ht]
\centering
\fig{ch07-chessboard}
\caption{The $64=65$ paradox. Left: an $8\times8$ square cut into two trapezia and two triangles along lines through lattice points. Right: the four pieces rearranged into a $5\times13$ rectangle. The ``diagonal'' of the rectangle is not a line: the edges of the pieces have slopes $\frac25$ and $\frac38$, and between them lies a thin parallelogram (shaded, drawn to scale) with vertices $(0,0)$, $(8,3)$, $(13,5)$, $(5,2)$. Its area is $\left|\det\begin{pmatrix}8&5\\3&2\end{pmatrix}\right|=|16-15|=1$: the missing square, smeared along the diagonal.}
\label{fig:chessboard}
\end{figure}

\portraits{cayley}{Arthur Cayley (1821--1895). Photograph by Herbert Rose Barraud, before 1883. Public domain, via Wikimedia Commons.}{hamilton}{William Rowan Hamilton (1805--1865). Portrait, artist unknown. Public domain, via Wikimedia Commons.}

\section{Vector spaces, from scratch}

\begin{reminder}[vector spaces, bases, dimension]
Let $\mathbb K$ be $\R$ or $\C$ (any field would do). A \emph{vector space} over $\mathbb K$ is a set $E$ with an addition and a multiplication by
scalars satisfying the usual rules ($\mathbb K^n$, polynomials, functions on a set, sequences are the examples to keep in mind); a \emph{subspace} is a
non-empty subset closed under both operations. Vectors $v_1,\dots,v_p$ are \emph{linearly independent} if $\lambda_1v_1+\dots+\lambda_pv_p=0$
forces all $\lambda_i=0$; they \emph{span} $E$ if every vector is such a combination; a \emph{basis} is a family that does both, so that every vector
has exactly one expression $\sum\lambda_iv_i$, the $\lambda_i$ being its \emph{coordinates}. The space $E$ is \emph{finite-dimensional} if it is spanned by finitely many vectors.
\end{reminder}

The word ``dimension'' is only meaningful because of the following theorem, which says that all bases of a finite-dimensional space have the same size.
It is the first non-trivial theorem of linear algebra, and everything else leans on it.

\begin{theoremx}[exchange lemma, and dimension]
Let $E$ be spanned by $v_1,\dots,v_n$.
\begin{enumerate}[label=(\alph*)]
\item \textbf{Steinitz's exchange lemma}: if $u_1,\dots,u_p$ are linearly independent in $E$, then $p\le n$.
\item $E$ has a basis, extracted from $v_1,\dots,v_n$; every linearly independent family extends to a basis; and all bases have the same number of elements, called $\dim E$.
\item A family of $\dim E$ vectors that is linearly independent, or that spans $E$, is a basis. Every subspace $F$ of $E$ is finite-dimensional with $\dim F\le\dim E$, and $\dim F=\dim E$ only if $F=E$.
\end{enumerate}
\end{theoremx}
\begin{proof}
(a) We show that the $u_i$ can replace the $v_j$ one at a time. Write $u_1=\sum_j\lambda_jv_j$; some $\lambda_j\ne0$ (since $u_1\ne0$), say $\lambda_1$ after renumbering, and then $v_1$ is a combination of $u_1,v_2,\dots,v_n$, so this family spans $E$. Suppose that after $k<p$ steps $u_1,\dots,u_k,v_{k+1},\dots,v_n$ spans $E$ (which requires $k\le n$). Write $u_{k+1}=\sum_{i\le k}\mu_iu_i+\sum_{j>k}\lambda_jv_j$. If all $\lambda_j$ were $0$, the $u_i$ would be dependent; so some $\lambda_j\ne0$, which forces $k<n$, and $v_j$ can be exchanged for $u_{k+1}$. After $p$ steps we have exchanged $p$ of the $v_j$, so $p\le n$.

(b) Among the subfamilies of $v_1,\dots,v_n$ that span $E$, take one of minimal size; it is linearly independent (a dependence relation would let us drop one vector and still span), hence a basis. If $u_1,\dots,u_p$ are independent, add to them the $v_j$ that are not combinations of the vectors already chosen, one by one; the family stays independent and ends up spanning $E$. Finally, if $B$ and $B'$ are bases with $|B|=b$, $|B'|=b'$, then (a) applied both ways gives $b\le b'$ and $b'\le b$.

(c) An independent family of $\dim E$ vectors extends to a basis, which has $\dim E$ elements, so nothing was added; a spanning family of $\dim E$ vectors contains a basis, again with nothing removed. For a subspace $F$, independent families in $F$ have at most $\dim E$ elements by (a), so a maximal one exists and is a basis of $F$ (a vector of $F$ outside its span could be added); if $\dim F=\dim E$ that basis is a basis of $E$ by the first claim, so $F=E$.
\end{proof}

\begin{reminder}[linear maps and matrices]
A map $u:E\to F$ is \emph{linear} if $u(\lambda x+\mu y)=\lambda u(x)+\mu u(y)$. Its \emph{kernel} $\ker u=u^{-1}(0)$ and \emph{image} $\operatorname{Im}u$ are subspaces, $u$ is injective iff $\ker u=\{0\}$, and the \emph{rank} of $u$ is $\dim\operatorname{Im}u$. Given bases $(e_j)$ of $E$ and $(f_i)$ of $F$, the \emph{matrix} of $u$ has in its $j$th column the coordinates of $u(e_j)$: $u(e_j)=\sum_ia_{ij}f_i$. Composition of maps is multiplication of matrices, $(AB)_{ik}=\sum_ja_{ij}b_{jk}$, which is why matrix multiplication is associative and not commutative. If $P$ is the matrix whose columns are the coordinates of a new basis $(e'_j)$ in the old one, then the coordinates $X$ and $X'$ of a vector in the two bases satisfy $X=PX'$, and the matrix of an endomorphism $u:E\to E$ in the new basis is $A'=P^{-1}AP$: $A$ and $A'$ are \emph{similar}. A statement about $u$ that does not depend on the basis must be invariant under $A\mapsto P^{-1}AP$; the trace, the determinant and the characteristic polynomial below are the basic invariants.
\end{reminder}

\begin{examplex}[the $2\times2$ case, worked]
In the plane, the matrix $R_\theta=\begin{pmatrix}\cos\theta&-\sin\theta\\ \sin\theta&\cos\theta\end{pmatrix}$ sends $e_1=(1,0)$ to $(\cos\theta,\sin\theta)$ and $e_2$ to $(-\sin\theta,\cos\theta)$: it is the rotation by $\theta$, and $R_\theta R_\varphi=R_{\theta+\varphi}$ by the addition formulas (Exercise~\ref{ex:twobytwo} makes this a statement about complex numbers). The matrix $\begin{pmatrix}1&1\\0&1\end{pmatrix}$ fixes $e_1$ and sends $e_2$ to $e_1+e_2$: a \emph{shear}, which slides each horizontal line along itself by its height; it preserves areas, as Figure~\ref{fig:parallelogram} shows, and it is the key to the area formula below. The matrix $\begin{pmatrix}a&b\\c&d\end{pmatrix}$ is invertible iff $ad-bc\ne0$, with inverse $\frac1{ad-bc}\begin{pmatrix}d&-b\\-c&a\end{pmatrix}$ (multiply and see); its kernel is $\{0\}$ iff its columns are not proportional iff $ad-bc\ne0$: the determinant, before it is defined, already decides invertibility.
\end{examplex}

\section{Determinants}

This section rebuilds the determinant from scratch, in a page. Let $\mathfrak S_n$ be the group of permutations of $\{1,\dots,n\}$.

\begin{lemma}[sign of a permutation]\label{lem:sign}
$\varepsilon(\sigma)=\prod_{i<j}\frac{\sigma(j)-\sigma(i)}{j-i}$ is $\pm1$ and satisfies $\varepsilon(\sigma\tau)=\varepsilon(\sigma)\varepsilon(\tau)$. Every transposition has sign $-1$, and every permutation is a product of transpositions.
\end{lemma}
\begin{proof}
The quotient $\frac{\sigma(j)-\sigma(i)}{j-i}$ does not change when $i$ and $j$ are swapped, so it depends only on the pair $\{i,j\}$, and $\varepsilon(\sigma)=\prod_{\text{pairs}}\frac{\sigma(j)-\sigma(i)}{j-i}$. As $\{i,j\}$ runs over the pairs, so does $\{\sigma(i),\sigma(j)\}$: numerator and denominator are the same product up to sign, so $\varepsilon(\sigma)=\pm1$. Next,
\[
\begin{aligned}\varepsilon(\sigma\tau)&=\prod_{\text{pairs}}\frac{\sigma(\tau(j))-\sigma(\tau(i))}{\tau(j)-\tau(i)}\\&\quad\cdot\prod_{\text{pairs}}\frac{\tau(j)-\tau(i)}{j-i}=\varepsilon(\sigma)\varepsilon(\tau),\end{aligned}
\]
the first product being $\varepsilon(\sigma)$ written over the pairs $\{\tau(i),\tau(j)\}$. For the transposition $(1\,2)$, only the pair $\{1,2\}$ gives a negative factor, so $\varepsilon=-1$; any transposition $(a\,b)$ equals $\pi(1\,2)\pi^{-1}$ for some $\pi$, so it has sign $\varepsilon(\pi)(-1)\varepsilon(\pi)^{-1}=-1$. Finally, by induction on $n$: if $\sigma(n)=k\ne n$, then $(k\,n)\sigma$ fixes $n$ and is a product of transpositions of $\{1,\dots,n-1\}$.
\end{proof}

Concretely, $\varepsilon(\sigma)=(-1)^{N}$ where $N$ is the number of \emph{inversions}, the pairs $i<j$ with $\sigma(i)>\sigma(j)$: each inversion contributes one negative factor. The permutation $(3,1,2)$ of $\{1,2,3\}$ (meaning $1\mapsto3$, $2\mapsto1$, $3\mapsto2$) has two inversions, so it is even; it is the $3$-cycle, a product of two transpositions.

\begin{theorem}[determinant]\label{thm:det}
\begin{enumerate}[label=(\alph*)]
\item There is a unique function $\det:M_n(\mathbb K)\to\mathbb K$ that is linear in each column, vanishes when two columns are equal (\emph{alternating}), and satisfies $\det I_n=1$; it is $\det A=\sum_{\sigma\in\mathfrak S_n}\varepsilon(\sigma)\prod_{j=1}^na_{\sigma(j),j}$. Every multilinear alternating $D$ satisfies $D(A)=D(I_n)\det A$.
\item $\det(AB)=\det A\det B$ and $\det A^{\mathsf T}=\det A$; in particular everything said about columns holds for rows.
\item $A$ is invertible if and only if $\det A\ne0$.
\item Adding to a column a multiple of another does not change the determinant, and the determinant of a triangular matrix is the product of its diagonal.
\end{enumerate}
\end{theorem}
\begin{proof}
(a) Let $D$ be multilinear and alternating. Expanding $D(x+y,x+y)=0$ in two columns shows that $D$ changes sign when two columns are swapped. Write column $j$ of $A$ as $\sum_ia_{ij}e_i$ and expand by multilinearity: $D(A)=\sum_{f}\prod_ja_{f(j),j}\,D(e_{f(1)},\dots,e_{f(n)})$ over all maps $f:\{1,\dots,n\}\to\{1,\dots,n\}$. Terms with $f$ non-injective vanish (two equal columns); for $f=\sigma$ a permutation, writing $\sigma$ as a product of $t$ transpositions and undoing them one at a time gives $D(e_{\sigma(1)},\dots,e_{\sigma(n)})=(-1)^tD(I_n)=\varepsilon(\sigma)D(I_n)$. This proves uniqueness and the formula $D(A)=D(I_n)\det A$.
Conversely, the formula is multilinear (each term contains exactly one entry of each column). If columns $k\ne l$ are equal, let $\tau=(k\,l)$; the map $\sigma\mapsto\sigma\tau$ pairs the permutations, and $\prod_ja_{\sigma\tau(j),j}=\prod_ja_{\sigma(j),j}$ (the columns $k,l$ are equal) while $\varepsilon(\sigma\tau)=-\varepsilon(\sigma)$: the terms cancel in pairs. And $\det I_n=1$ (only $\sigma=\mathrm{id}$ contributes).

(b) For fixed $A$, $X\mapsto\det(AX)$ is multilinear and alternating in the columns of $X$ (column $j$ of $AX$ is $A$ times column $j$ of $X$), so by (a) it equals $\det(A)\det X$. For the transpose, substitute $\sigma\to\sigma^{-1}$ in the formula, using $\varepsilon(\sigma^{-1})=\varepsilon(\sigma)$ and $\prod_ja_{\sigma(j),j}=\prod_ia_{i,\sigma^{-1}(i)}$.

(c) If $AA^{-1}=I_n$ then $\det A\det A^{-1}=1$. If $A$ is not invertible, its columns are dependent, say $C_k=\sum_{j\ne k}\lambda_jC_j$; by linearity in column $k$, $\det A$ is a sum of determinants with two equal columns, so $0$.

(d) The first claim is linearity in that column plus alternation. For an upper triangular matrix, a non-zero term of the formula needs $\sigma(j)\le j$ for all $j$, which forces $\sigma=\mathrm{id}$ (by induction from $j=1$). Lower triangular: transpose.
\end{proof}

The formula has $n!$ terms, which makes it useless for computing beyond $n=3$; what one uses instead is (d), to reduce a matrix to triangular form by
row operations in about $n^3$ operations, or the following expansion, which reduces an $n\times n$ determinant to $n$ determinants of size $n-1$.

\begin{propositionx}[Laplace expansion]
Let $A_{ij}$ be the $(n-1)\times(n-1)$ matrix obtained from $A$ by deleting row $i$ and column $j$, and $c_{ij}=(-1)^{i+j}\det A_{ij}$ the \emph{cofactor}. Then for every row $i$ and every column $j$,
\[
\det A=\sum_{k=1}^na_{ik}c_{ik}=\sum_{k=1}^na_{kj}c_{kj}.
\]
\end{propositionx}
\begin{proof}
Expansion along column $j$: by linearity in that column, $\det A=\sum_ka_{kj}\det A^{(k)}$, where $A^{(k)}$ is $A$ with column $j$ replaced by $e_k$. Move column $j$ to the first position by $j-1$ swaps of adjacent columns, then row $k$ to the first position by $k-1$ swaps of adjacent rows: the sign changes $(-1)^{j-1+k-1}=(-1)^{k+j}$ times, and the matrix now has first column $e_1$ and, below and to the right of the corner, the matrix $A_{kj}$. In the formula of Theorem~\ref{thm:det}(a), the only non-zero terms have $\sigma(1)=1$, and they form the determinant of $A_{kj}$. So $\det A^{(k)}=(-1)^{k+j}\det A_{kj}$. Rows: transpose.
\end{proof}

\begin{examplex}
Along the first row,
\[
\begin{aligned}
\det\begin{pmatrix}2&1&3\\0&4&1\\1&2&5\end{pmatrix}
&=2\det\begin{pmatrix}4&1\\2&5\end{pmatrix}-\det\begin{pmatrix}0&1\\1&5\end{pmatrix}
+3\det\begin{pmatrix}0&4\\1&2\end{pmatrix}\\
&=2\cdot18-(-1)+3\cdot(-4)=25.
\end{aligned}
\]
Along the first column it is $2\cdot18-0+1\cdot(1-12)=25$ too, as it must be. The signs $(-1)^{i+j}$ form a chessboard pattern starting with $+$ in the corner; the row or column with the most zeros is the one to expand along.
\end{examplex}

\begin{proposition}[rank--nullity]\label{prop:ranknullity}
For $u:E\to F$ linear with $\dim E<\infty$, $\dim\ker u+\operatorname{rk}u=\dim E$.
\end{proposition}
\begin{proof}
Take a basis $(k_1,\dots,k_p)$ of $\ker u$ and extend it to a basis $(k_1,\dots,k_p,e_1,\dots,e_r)$ of $E$ (possible by the exchange theorem). The vectors $u(e_i)$ span $\operatorname{Im}u$ (the $u(k_i)$ vanish), and they are independent: $\sum\lambda_iu(e_i)=0$ puts $\sum\lambda_ie_i$ in $\ker u$, so it is a combination of the $k_i$, and all coefficients vanish. So $\operatorname{rk}u=r=\dim E-p$.
\end{proof}

For a matrix, the rank is the dimension of the span of its columns, and rank--nullity says that an $n\times n$ matrix is invertible iff its kernel is $\{0\}$ iff its columns are independent iff its rank is $n$; Theorem~\ref{thm:det}(c) adds ``iff $\det A\ne0$''. The rank of the rows equals the rank of the columns (Exercise~\ref{ex:tracerank}).

\begin{think}
Compute $\det\begin{pmatrix}1&a\\1&b\end{pmatrix}$, then guess $\det\begin{pmatrix}1&a&a^2\\1&b&b^2\\1&c&c^2\end{pmatrix}$ without expanding: it is a polynomial in $a,b,c$ of total degree $3$ that vanishes when two of the variables are equal.
\end{think}

\begin{proposition}[Vandermonde]\label{prop:vdm}
For $x_1,\dots,x_n\in\mathbb K$,
\[
\begin{aligned}V(x_1,\dots,x_n)&=\det\left(x_i^{\,j-1}\right)_{1\le i,j\le n}\eqbr=\prod_{1\le i<j\le n}(x_j-x_i).\end{aligned}
\]
\end{proposition}
\begin{proof}
Induction on $n$; $n=1$ is clear. If two of the $x_i$ are equal both sides vanish. Otherwise, replacing $x_n$ by a variable $X$ and expanding along the last row (Laplace),
$P(X)=V(x_1,\dots,x_{n-1},X)$ is a polynomial of degree at most $n-1$ whose coefficient of $X^{n-1}$ is the cofactor of the corner entry, $V(x_1,\dots,x_{n-1})$. It vanishes at $x_1,\dots,x_{n-1}$ (two equal rows),
so $P=V(x_1,\dots,x_{n-1})\prod_{i<n}(X-x_i)$ by Proposition~\ref{prop:rootsfactor}. Evaluate at $X=x_n$ and apply the induction hypothesis.
\end{proof}

The Vandermonde determinant is what makes Lagrange interpolation (Theorem~\ref{thm:lagrange}) a statement about an invertible matrix: the system ``find a polynomial of degree $<n$ with prescribed values at $n$ distinct points'' has the matrix $(x_i^{j-1})$, invertible exactly when the points are distinct. (The circulant determinant of Exercise~\ref{ex:circulant} is another closed form, obtained by diagonalising.)

\section{Reduction}

A matrix is a linear map written in a basis, and a good choice of basis can make it almost transparent: diagonal, if we are lucky. The numbers on the
diagonal are then the \emph{eigenvalues}, the directions of the basis the \emph{eigenvectors}, and the map is a stretch by $\lambda_i$ along the
$i$th direction. The programme of this section is to decide when such a basis exists, and to find it. Let $A\in M_n(\mathbb K)$ and $\chi_A(X)=\det(XI_n-A)$.

\begin{reminder}[eigenvalues, eigenvectors, similarity]
A scalar $\lambda\in\mathbb K$ is an \emph{eigenvalue} of $A$ if $Av=\lambda v$ for some $v\ne0$, an \emph{eigenvector}; the \emph{eigenspace} $\ker(A-\lambda I_n)$ collects them, with $0$. Two matrices $A$ and $B=P^{-1}AP$ with $P$ invertible are \emph{similar}: they are the same linear map written in two bases, $P$ holding the coordinates of the new basis vectors. $A$ is \emph{diagonalisable} if it is similar to a diagonal matrix, that is, if $\mathbb K^n$ has a basis of eigenvectors of $A$. A \emph{polynomial in} $A$ is $P(A)=\sum p_kA^k$ for $P=\sum p_kX^k$, with $A^0=I_n$; $P\mapsto P(A)$ respects sums and products, since the powers of $A$ commute with each other.
\end{reminder}

\begin{proposition}[characteristic polynomial]\label{prop:charpoly}
$\chi_A=X^n-\operatorname{tr}(A)X^{n-1}+\dots+(-1)^n\det A$, and the eigenvalues of $A$ are its roots in $\mathbb K$. Over $\C$, $\operatorname{tr}A$ and $\det A$ are the sum and the product of the eigenvalues, counted with multiplicity as roots of $\chi_A$. Similar matrices have the same $\chi$.
\end{proposition}
\begin{proof}
In the formula for $\det(XI_n-A)$, a permutation $\sigma\ne\mathrm{id}$ moves at least two indices, so its term contains at most $n-2$ diagonal entries $X-a_{jj}$ and has degree $\le n-2$. So the coefficients of $X^n$ and $X^{n-1}$ come from $\prod_j(X-a_{jj})$: they are $1$ and $-\operatorname{tr}A$. The constant term is $\chi_A(0)=\det(-A)=(-1)^n\det A$. Next, $\lambda$ is an eigenvalue iff $\ker(\lambda I_n-A)\ne\{0\}$ iff $\det(\lambda I_n-A)=0$ (Theorem~\ref{thm:det}(c)). Over $\C$, $\chi_A$ splits (Theorem~\ref{thm:dalembert}), and Vi\`ete gives the sum and product of the roots. Finally $\det(XI-P^{-1}AP)=\det\left(P^{-1}(XI-A)P\right)=\chi_A$.
\end{proof}

\begin{examplex}[a $2\times2$ matrix, diagonalised by hand]
For $A=\begin{pmatrix}2&1\\1&2\end{pmatrix}$, $\chi_A=(X-2)^2-1=(X-1)(X-3)$. For $\lambda=1$, $A-I=\begin{pmatrix}1&1\\1&1\end{pmatrix}$ kills $u_1=(1,-1)$; for $\lambda=3$, $A-3I=\begin{pmatrix}-1&1\\1&-1\end{pmatrix}$ kills $u_2=(1,1)$. With $P=\begin{pmatrix}1&1\\-1&1\end{pmatrix}$ (the eigenvectors as columns), $AP=P\operatorname{diag}(1,3)$, so $A=PDP^{-1}$ and $A^n=PD^nP^{-1}=\frac12\begin{pmatrix}3^n+1&3^n-1\\3^n-1&3^n+1\end{pmatrix}$. Check $n=1$. The trace $4=1+3$ and the determinant $3=1\cdot3$ agree with Proposition~\ref{prop:charpoly}. On the other hand $N=\begin{pmatrix}0&1\\0&0\end{pmatrix}$ has $\chi_N=X^2$, a single eigenvalue $0$ whose eigenspace $\ker N=\mathbb Ke_1$ is a line: no basis of eigenvectors, so $N$ is not diagonalisable. The failure is not a lack of eigenvalues but a lack of eigenvectors.
\end{examplex}

\begin{think}
A projection satisfies $A^2=A$, and a reflection satisfies $A^2=I$. Both are diagonalisable. Is every matrix with $A^2=0$ diagonalisable? What distinguishes $X^2-X$ and $X^2-1$ from $X^2$?
\end{think}

\begin{theorem}[diagonalisability]\label{thm:diag}
$A$ is diagonalisable if and only if $P(A)=0$ for some polynomial $P=\prod_{i=1}^r(X-\lambda_i)$ with distinct $\lambda_i\in\mathbb K$. In that case $\mathbb K^n=\bigoplus_i\ker(A-\lambda_iI_n)$.
\end{theorem}
\begin{proof}
If $A=QDQ^{-1}$ with $D$ diagonal, let $\lambda_1,\dots,\lambda_r$ be the distinct diagonal entries: $\prod(D-\lambda_iI)=0$ (each diagonal entry is killed by one factor), hence $\prod(A-\lambda_iI)=Q\prod(D-\lambda_iI)Q^{-1}=0$.
Conversely let $L_i$ be the Lagrange polynomials on the points $\lambda_1,\dots,\lambda_r$ (Theorem~\ref{thm:lagrange}). Interpolating the constant $1$ gives $\sum_iL_i=1$, so every $v\in\mathbb K^n$ is $v=\sum_iL_i(A)v$. And $(X-\lambda_i)L_i$ is a constant multiple of $P$, so $(A-\lambda_iI)L_i(A)v=0$: $L_i(A)v\in\ker(A-\lambda_iI)$. So the eigenspaces span $\mathbb K^n$. The sum is direct: if $\sum v_i=0$ with $Av_i=\lambda_iv_i$, apply $L_j(A)$, which acts on $v_i$ as multiplication by $L_j(\lambda_i)=\delta_{ij}$: $v_j=0$. A basis adapted to the decomposition diagonalises $A$.
\end{proof}

The theorem answers the question in the box: $A^2=0$ gives the annihilating polynomial $X^2$, with a double root, which proves nothing; and indeed $N$ above is a counterexample. But $A^2=A$ and $A^2=I$ give $X(X-1)$ and $(X-1)(X+1)$, split with simple roots, so projections and reflections are diagonalisable (over any field in which $1\ne-1$, for reflections). Two useful consequences: a matrix whose characteristic polynomial has $n$ distinct roots in $\mathbb K$ is diagonalisable, since $\chi_A$ then has the required form and $\chi_A(A)=0$ by Theorem~\ref{thm:cayleyhamilton} below (for this special case one can avoid Cayley--Hamilton: eigenvectors for distinct eigenvalues are independent, by the same $L_j(A)$ trick, and $n$ of them form a basis); and a diagonalisable matrix stays diagonalisable on every stable subspace, since the same polynomial kills the restriction (Exercise~\ref{ex:codiag} uses this).

\begin{remark}
Cayley--Hamilton, the statement that $\chi_A(A)=0$, is proved below as Theorem~\ref{thm:cayleyhamilton}, by a trick with the adjugate. Volume~2 proves it a second time, as a consequence of the structure of modules over $\mathbb K[X]$, and develops the minimal polynomial, which is not needed here.
\end{remark}

\begin{think}
An invertible $2\times2$ symmetric matrix maps the unit circle to an ellipse (Figure~\ref{fig:ellipse}). An ellipse has two perpendicular axes. What should they be for the matrix, and what are their lengths?
\end{think}

\begin{theorem}[spectral theorem]\label{thm:spectral}
A real symmetric matrix $S$ has real eigenvalues and an orthonormal basis of eigenvectors: $S=PDP^{\mathsf T}$ with $P$ orthogonal and $D$ diagonal.
\end{theorem}
\begin{proof}
Let $\lambda\in\C$ be a root of $\chi_S$ and $z\in\C^n\setminus\{0\}$ with $Sz=\lambda z$. Then $\bar z^{\mathsf T}Sz=\lambda\|z\|^2$, and this number is real: its conjugate is $z^{\mathsf T}S\bar z$, which is the same scalar transposed ($S^{\mathsf T}=S$). So $\lambda\in\R$, and $\ker(S-\lambda I)$ contains a real unit vector $u$.
Induction on $n$: the hyperplane $H=u^\perp$ is stable under $S$, since $x\perp u$ implies $\langle Sx,u\rangle=\langle x,Su\rangle=\lambda\langle x,u\rangle=0$. The restriction of $S$ to $H$ is symmetric for the induced inner product, so $H$ has an orthonormal basis of eigenvectors; add $u$.
\end{proof}

\begin{figure}[ht]
\centering
\fig{ch07-ellipse}
\caption{Left: $S=\begin{pmatrix}2&1\\1&2\end{pmatrix}$ maps the unit circle (grey) to an ellipse (blue). Each red arrow joins a point $x$ of the circle to $Sx$. The axes of the ellipse are the orthonormal eigenvectors $u_1=\frac1{\sqrt2}(-1,1)$ and $u_2=\frac1{\sqrt2}(1,1)$, stretched by the eigenvalues $1$ and $3$, the extreme values of $x^{\mathsf T}Sx$ on the circle (Exercise~\ref{ex:rayleigh}). Right: $\det H_n$ for the Hilbert matrix of Exercise~\ref{ex:hilbert}, on a logarithmic scale. It is positive but falls from $\frac1{12}$ for $n=2$ to about $2.6\times10^{-78}$ for $n=12$.}
\label{fig:ellipse}
\end{figure}

The spectral theorem is the reason the ellipse of Figure~\ref{fig:ellipse} has perpendicular axes, and the reason the physicists' ``principal axes'' of a rotating body exist: Cauchy's 1829 proof was about exactly that, a quadratic form in three variables and its axes, and the matrix was not yet invented. In Volume~3 the same theorem diagonalises covariance matrices, which is what ``principal component analysis'' means.

\begin{example}[linear recurrences]
If $u_{n+2}=au_{n+1}+bu_n$, then $\binom{u_{n+1}}{u_n}=M^n\binom{u_1}{u_0}$ with $M=\begin{pmatrix}a&b\\1&0\end{pmatrix}$, whose characteristic polynomial is $X^2-aX-b$.
When $b\ne0$ and it has two distinct roots $r_1,r_2$, $M$ is diagonalisable by Theorem~\ref{thm:diag}, and $u_n=\alpha r_1^n+\beta r_2^n$ for constants $\alpha,\beta$ fixed by $u_0,u_1$.
\end{example}

If $b=0$, handle the recurrence directly: $u_1$ is arbitrary and $u_n=a^{n-1}u_1$ for $n\ge2$, while $u_0$ is arbitrary; when $a=0$, all $u_n$ with $n\ge2$ are zero. This avoids treating $0^0$ as an unexplained initial condition. The double-root case is the other case the example leaves out: for $u_{n+2}=2u_{n+1}-u_n$ the matrix $\begin{pmatrix}2&-1\\1&0\end{pmatrix}$ has $\chi=(X-1)^2$ and is not diagonalisable (its only eigenspace is the line spanned by $(1,1)$), and the solutions are $u_n=\alpha+\beta n$: the factor $n$ is the trace that non-diagonalisability leaves behind. Volume~2 explains it with the Jordan form.

\section{Determinants as volumes}

Back to the chessboard. The gap in Figure~\ref{fig:chessboard} is a parallelogram, and its area was a $2\times2$ determinant. This is not a coincidence of the puzzle but the geometric meaning of determinants: $|\det A|$ is the volume of the parallelepiped spanned by the columns of $A$, and the sign records the orientation. To make sense of ``volume'' in $\R^n$ without any theory of measure, we use the inner product $\langle x,y\rangle=\sum x_iy_i$ and the norm $\|x\|=\sqrt{\langle x,x\rangle}$, and we define the volume as the engineers do, base times height, by induction. The tool is orthogonal projection.

\begin{figure}[ht]
\centering
\fig{ch07-parallelogram}
\caption{Left: the parallelogram spanned by the columns $v_1=(a,c)$ and $v_2=(b,d)$ of a matrix; its area is $|ad-bc|$. Middle: a shear, adding a multiple of $v_1$ to $v_2$, slides the top side along itself and changes neither the area nor the determinant (Theorem~\ref{thm:det}(d)). Right: after shearing $v_2$ into $w_2\perp v_1$, the parallelogram is a rectangle of area $\|v_1\|\cdot\|w_2\|$, and $\|w_2\|$ is the distance from $v_2$ to the line $\R v_1$.}
\label{fig:parallelogram}
\end{figure}

\begin{theorem}[Gram--Schmidt and orthogonal projection]\label{thm:gramschmidt}
Let $F$ be a subspace of $\R^n$ with basis $v_1,\dots,v_p$.
\begin{enumerate}[label=(\alph*)]
\item (\textbf{Gram--Schmidt}) There is an orthonormal family $e_1,\dots,e_p$ with $\operatorname{Span}(e_1,\dots,e_k)=\operatorname{Span}(v_1,\dots,v_k)$ for each $k$, given by $w_k=v_k-\sum_{i<k}\langle v_k,e_i\rangle e_i$ and $e_k=w_k/\|w_k\|$.
\item (\textbf{projection}) Every $x\in\R^n$ is uniquely $x=p(x)+r$ with $p(x)\in F$ and $r\perp F$; explicitly $p(x)=\sum_{i=1}^p\langle x,e_i\rangle e_i$. The map $p$ is linear with $p^2=p$, and $p(x)$ is the point of $F$ closest to $x$:
\[
\|x-p(x)\|<\|x-y\|\quad\text{for every }y\in F,\ y\ne p(x).
\]
We call $d(x,F)=\|x-p(x)\|$ the \emph{distance} from $x$ to $F$.
\end{enumerate}
\end{theorem}
\begin{proof}
(a) Induction on $k$. Suppose $e_1,\dots,e_{k-1}$ orthonormal with the same span as $v_1,\dots,v_{k-1}$. Then $w_k$ is $v_k$ minus a vector of that span, so $w_k\ne0$ ($v_k$ is not in the span of the earlier $v_i$, by independence), and for $j<k$,
\[
\langle w_k,e_j\rangle=\langle v_k,e_j\rangle-\sum_{i<k}\langle v_k,e_i\rangle\langle e_i,e_j\rangle=\langle v_k,e_j\rangle-\langle v_k,e_j\rangle=0.
\]
Thus $e_k=w_k/\|w_k\|$ extends the orthonormal family. Moreover,
\[
\operatorname{Span}(e_1,\dots,e_k)\subseteq\operatorname{Span}(v_1,\dots,v_k).
\]
Both spaces have dimension $k$, so they are equal.

(b) Put $p(x)=\sum_{i\le p}\langle x,e_i\rangle e_i\in F$ and $r=x-p(x)$. For each $j$, $\langle r,e_j\rangle=\langle x,e_j\rangle-\langle x,e_j\rangle=0$, so $r\perp F$. Uniqueness: if $x=f+r=f'+r'$ with $f,f'\in F$ and $r,r'\perp F$, then $f-f'=r'-r$ lies in $F$ and is orthogonal to $F$, so it is orthogonal to itself: $f=f'$. Linearity and $p^2=p$ are read off the formula ($p(e_i)=e_i$). For $y\in F$, $x-y=(x-p(x))+(p(x)-y)$ with the two terms orthogonal, so by Pythagoras
\[
\|x-y\|^2=\|x-p(x)\|^2+\|p(x)-y\|^2\ge\|x-p(x)\|^2,
\]
with equality only for $y=p(x)$.
\end{proof}

Gram--Schmidt is both an algorithm and a proof that every subspace has an orthonormal basis, which the spectral theorem used. The closest-point property is the start of least squares, the method by which Gauss located the lost asteroid Ceres in 1801: ``the best approximation to the data within the subspace of what the model allows'' is an orthogonal projection.

\begin{definitionx}[volume]
Let $v_1$, $\dots$, $v_p$ belong to $\R^n$. Define the successive heights by
\[
\begin{gathered}
d_1=\|v_1\|,\\
d_k=d\left(v_k,\operatorname{Span}(v_1,\dots,v_{k-1})\right)\quad(k\ge2).
\end{gathered}
\] The \emph{$p$-dimensional volume} of the parallelepiped $\{\sum t_iv_i:0\le t_i\le1\}$ is $\operatorname{vol}(v_1,\dots,v_p)=d_1d_2\cdots d_p$: base times height, repeated.
\end{definitionx}

\begin{theorem}[Gram determinant]\label{thm:gram}
Let $v_1,\dots,v_p\in\R^n$ and $G=\left(\langle v_i,v_j\rangle\right)_{1\le i,j\le p}$, the \emph{Gram matrix}. Then
\[
\det G=\operatorname{vol}(v_1,\dots,v_p)^2 ,
\]
so $\det G\ge0$, with equality if and only if $v_1,\dots,v_p$ are linearly dependent. If $p=n$ and $A$ is the matrix with columns $v_1,\dots,v_n$, then $G=A^{\mathsf T}A$ and
\[
|\det A|=\operatorname{vol}(v_1,\dots,v_n).
\]
\end{theorem}
\begin{proof}
If the $v_i$ are dependent, say $\sum c_iv_i=0$ with $c\ne0$, then $Gc=0$ (its $i$th entry is $\langle v_i,\sum c_jv_j\rangle$), so $\det G=0$; and some $d_k=0$, since the first $v_k$ that is a combination of the earlier ones lies in their span, so the volume is $0$ too. Suppose the $v_i$ independent and run Gram--Schmidt: $w_k=v_k-\sum_{i<k}\langle v_k,e_i\rangle e_i$, where the subtracted vector is $p_k(v_k)$, the projection on $\operatorname{Span}(v_1,\dots,v_{k-1})$; so $\|w_k\|=d_k$ by definition of the distance. Replacing $v_k$ by $w_k$ replaces, in $G$, row $k$ by row $k$ minus a combination of the earlier rows, and column $k$ likewise; by Theorem~\ref{thm:det}(d), applied to rows and to columns, $\det G$ does not change. Do this for $k=2,\dots,p$ in order (each step only uses the $w_i$, $i<k$, whose span is that of the $v_i$, $i<k$). At the end $G$ has become the Gram matrix of $w_1,\dots,w_p$, which are pairwise orthogonal: it is $\operatorname{diag}(\|w_1\|^2,\dots,\|w_p\|^2)$, with determinant $\prod d_k^2=\operatorname{vol}^2$. For $p=n$, the $(i,j)$ entry of $A^{\mathsf T}A$ is $\langle v_i,v_j\rangle$, so $\det G=\det(A^{\mathsf T})\det A=(\det A)^2$ by Theorem~\ref{thm:det}(b).
\end{proof}

\begin{figure}[ht]
\centering
\fig{ch07-gramschmidt}
\caption{Gram--Schmidt in the plane, for $v_1=(3,1)$ and $v_2=(1,2)$. The projection of $v_2$ on the line $\R v_1$ is $\langle v_2,e_1\rangle e_1=\frac12v_1$; the remainder $w_2=v_2-\frac12v_1=(-\frac12,\frac32)$ is perpendicular to $v_1$, and $\|w_2\|=\sqrt{10}/2$ is the distance from $v_2$ to the line. The area of the parallelogram is $\|v_1\|\,\|w_2\|=\sqrt{10}\cdot\sqrt{10}/2=5=\det\begin{pmatrix}3&1\\1&2\end{pmatrix}$, and the Gram determinant is $\det\begin{pmatrix}10&5\\5&5\end{pmatrix}=25=5^2$.}
\label{fig:gramschmidt}
\end{figure}

Three things are hidden in this theorem. For $p=2$ it says $\|u\|^2\|v\|^2-\langle u,v\rangle^2\ge0$: the Cauchy--Schwarz inequality of Chapter~4, with its equality case, now as the statement that a parallelogram has non-negative area. For $p=n$ it says that the volume of the image of the unit cube under $A$ is $|\det A|$, and since $\det(AB)=\det A\det B$, the volume of the image of \emph{any} parallelepiped is multiplied by $|\det A|$: this is the meaning of the Jacobian in the change-of-variables formula of Volume~4. And it explains the Hilbert matrix of the epigraph: $H_n$ is the Gram matrix of $1,x,\dots,x^{n-1}$ for an inner product on polynomials (Exercise~\ref{ex:hilbert}), so $\det H_n$ is a squared volume, positive; it is tiny because the monomials become nearly linearly dependent in this inner product, so the parallelepiped is almost flat. Hadamard's inequality (Exercise~\ref{ex:hadamard}) is the same picture: a parallelepiped with given edge lengths has the largest volume when it is a box.

\section{The adjugate and Cayley--Hamilton}

The Laplace expansion computes $\det A$ from the cofactors $c_{ij}$. Expanding along the wrong row, with the cofactors of one row and the entries of another, gives $0$, and this harmless remark produces the inverse of a matrix in closed form.

\begin{proposition}[adjugate, Cramer's rule]\label{prop:adjugate}
Let $\operatorname{adj}A=(c_{ji})_{i,j}$ be the transpose of the matrix of cofactors. Then
\[
A\operatorname{adj}A=\operatorname{adj}A\;A=(\det A)\,I_n ,
\]
so if $\det A\ne0$, $A^{-1}=\frac1{\det A}\operatorname{adj}A$. In that case the solution of $Ax=b$ is given by \textbf{Cramer's rule}: $x_j=\det A_j/\det A$, where $A_j$ is $A$ with column $j$ replaced by $b$.
\end{proposition}
\begin{proof}
The $(i,k)$ entry of $A\operatorname{adj}A$ is $\sum_ja_{ij}c_{kj}$. For $i=k$ this is the Laplace expansion of $\det A$ along row $i$. For $i\ne k$, it is the expansion along row $k$ of the matrix $A'$ obtained from $A$ by replacing row $k$ with a copy of row $i$: the cofactors of row $k$ do not involve row $k$, so they are the same for $A$ and $A'$, and the entries of row $k$ of $A'$ are the $a_{ij}$. But $A'$ has two equal rows, so $\det A'=0$. The same with columns gives $\operatorname{adj}A\;A=(\det A)I_n$. For Cramer, $x=A^{-1}b=\frac1{\det A}\operatorname{adj}A\,b$ has $j$th entry $\frac1{\det A}\sum_kc_{kj}b_k$, and $\sum_kc_{kj}b_k$ is the expansion of $\det A_j$ along its column $j$, whose entries are the $b_k$ and whose cofactors are those of $A$.
\end{proof}

For $n=2$ this is the formula of the worked example: $\operatorname{adj}\begin{pmatrix}a&b\\c&d\end{pmatrix}=\begin{pmatrix}d&-b\\-c&a\end{pmatrix}$. For large $n$, Cramer's rule is a theoretical statement, not an algorithm: it costs $n+1$ determinants, where Gaussian elimination solves the system directly in about $n^3/3$ multiplications. What it is good for is \emph{structure}: it shows that the entries of $A^{-1}$ are polynomials in the entries of $A$ divided by $\det A$, which is how Exercise~\ref{ex:intinverse} decides which integer matrices have integer inverses, and it is the key to the following theorem, stated by Cayley in 1858 and proved in general by Frobenius in 1878.

\begin{theorem}[Cayley--Hamilton]\label{thm:cayleyhamilton}
Every square matrix satisfies its own characteristic equation: $\chi_A(A)=0$.
\end{theorem}

Before the proof, a warning about a proof that is too good to be true. ``Since $\chi_A(X)=\det(XI-A)$, substitute $X=A$: $\chi_A(A)=\det(AI-A)=\det0=0$.'' This substitutes a matrix for $X$ in two different ways at once: inside $\det(XI-A)$, $X$ is a scalar placed on the diagonal, and $AI-A$ would have to mean ``$A$ on each diagonal entry'', which is a block matrix, not $0$. The correct proof uses the adjugate of $XI-A$, a matrix whose entries are polynomials in $X$, and compares coefficients.

\begin{proof}
Write $\chi_A=X^n+c_{n-1}X^{n-1}+\dots+c_1X+c_0$. The cofactors of $XI_n-A$ are determinants of $(n-1)\times(n-1)$ matrices with entries in $\mathbb K[X]$ of degree $\le1$, so they are polynomials of degree at most $n-1$. Therefore
\[
\operatorname{adj}(XI_n-A)=B_0+XB_1+\dots+X^{n-1}B_{n-1}
\]
for some matrices $B_k\in M_n(\mathbb K)$ (collect the coefficient of $X^k$ in each entry). Proposition~\ref{prop:adjugate}, which is a polynomial identity valid over any commutative ring, here $\mathbb K[X]$, gives
\[
(XI_n-A)\left(B_0+XB_1+\dots+X^{n-1}B_{n-1}\right)=\chi_A(X)\,I_n.
\]
Compare the coefficients of each power of $X$ (two matrices of polynomials are equal iff all coefficients match):
\[
\begin{aligned}
X^n:&\quad B_{n-1}=I_n,\\
X^k\ (1\le k\le n-1):&\quad B_{k-1}-AB_k=c_kI_n,\\
X^0:&\quad -AB_0=c_0I_n.
\end{aligned}
\]
Multiply the equation of $X^k$ on the left by $A^k$ and add them all up: the left side is
\[
A^nB_{n-1}+\sum_{k=1}^{n-1}\left(A^kB_{k-1}-A^{k+1}B_k\right)-AB_0 ,
\]
a telescoping sum whose terms cancel in pairs ($A^kB_{k-1}$ cancels $-A^{k}B_{k-1}$ from the next equation down), leaving $0$. The right side is $A^n+\sum_{k<n}c_kA^k=\chi_A(A)$.
\end{proof}

Note what the telescoping used: only that $A^k$ multiplied from the left commutes with nothing in particular. We never needed the $B_k$ to commute with $A$ (they do, in fact, but proving it would be an extra step).

\begin{examplex}[Cayley--Hamilton at work]
For $n=2$: $A^2-(\operatorname{tr}A)A+(\det A)I_2=0$, which the reader should verify by hand once in their life. For the Fibonacci matrix $M=\begin{pmatrix}1&1\\1&0\end{pmatrix}$, $M^2=M+I$, so every power of $M$ is $aM+bI$, and Exercise~\ref{ex:fibch} turns this into identities about Fibonacci numbers. In general, $\chi_A(A)=0$ expresses $A^n$ as a combination of $I,A,\dots,A^{n-1}$, so \emph{every} power $A^m$, and every polynomial in $A$, lies in the span of $I,A,\dots,A^{n-1}$: the algebra generated by a matrix has dimension at most $n$. And when $c_0=(-1)^n\det A\ne0$, dividing $\chi_A(A)=0$ by $c_0$ and factoring $A$ out of the rest gives
\[
A^{-1}=-\frac1{c_0}\left(A^{n-1}+c_{n-1}A^{n-2}+\dots+c_1I_n\right),
\]
the inverse as a polynomial in the matrix: Exercise~\ref{ex:annul} was a special case.
\end{examplex}

\begin{beauty}
Cassini's identity comes from $\det(M^n)=(\det M)^n$: an arithmetic fact read off a determinant. The spectral theorem diagonalises the symmetric covariance matrices of
Volume~3. General Markov chains and PageRank in Exercise~\ref{ex:stochastic} use eigenvectors of matrices that need not be symmetric. Simultaneous diagonalisation (Exercise~\ref{ex:codiag}) is the
seed of the representation theory of Volume~2, where a finite abelian group of complex matrices is diagonalised at once by characters.
The Hilbert matrix is positive definite, yet it is ill-conditioned in the monomial basis, so ordinary floating-point inversion can lose many digits (a small determinant alone does not diagnose ill-conditioning): see Exercise~\ref{ex:cauchydet} for its exact determinant.
\end{beauty}

\portraits{sylvester}{James Joseph Sylvester (1814--1897). Photograph, nineteenth century. Public domain, via Wikimedia Commons.}{vandermonde}{Alexandre-Th\'eophile Vandermonde (1735--1796). Portrait, eighteenth century. Public domain, via Wikimedia Commons.}

\section{Workshop: least squares, with the geometry left in}
Three measurements rarely lie on a perfect straight line. A matrix does not have to bully the data into agreeing: it can choose the closest line and report the discrepancy.

\begin{examplex}[three points and an honest line]
Fit $y=a+bx$ to $(0,1),(1,2),(2,2)$ by minimising the sum of squared vertical errors. Write
\[
A=\begin{pmatrix}1&0\\1&1\\1&2\end{pmatrix},\qquad
Y=\begin{pmatrix}1\\2\\2\end{pmatrix}.
\]
The columns are independent. The closest vector $A\theta$ to $Y$ in their span satisfies $A^{\mathsf T}(Y-A\theta)=0$ by orthogonal projection, so
\[
\begin{pmatrix}3&3\\3&5\end{pmatrix}\binom ab=\binom56.
\]
The right side is $A^{\mathsf T}Y=(1+2+2,\,0\cdot1+1\cdot2+2\cdot2)^{\mathsf T}=(5,6)^{\mathsf T}$. Subtracting the first equation from the second gives $2b=1$. Hence $b=1/2$ and $a=7/6$. The residual vector is $(-1/6,1/3,-1/6)^{\mathsf T}$, orthogonal to both columns of $A$. Its squared length is $1/6$.
\end{examplex}

\begin{figure}[ht]\centering\fig{review07-fit}
\caption{The least-squares line $y=7/6+x/2$ and the three observations. Dashed vertical segments show the residuals. We minimise their squared lengths, not the perpendicular distances in this drawing.}\end{figure}

\subsection*{Why is this actually a minimum?}
\textbf{Mission.} Prove uniqueness without differentiating a quadratic in two variables.

\textbf{Hint.} Orthogonality turns a squared norm of a sum into a sum of squared norms.

\textbf{Solution.} Let $\widehat\theta=(7/6,1/2)^{\mathsf T}$ and $r=Y-A\widehat\theta$. For every $h\in\R^2$, $r\perp Ah$, so
\[
\|Y-A(\widehat\theta+h)\|^2=\|r\|^2+\|Ah\|^2\ge\|r\|^2.
\]
Equality forces $Ah=0$, and column independence gives $h=0$. This proves both minimality and uniqueness. For dependent columns, the closest fitted \emph{vector} is still unique, but its coefficient vector need not be.

\begin{reminder}[small determinant, different question]
For $0<\varepsilon<1$, the matrix $\varepsilon I$ has tiny determinant $\varepsilon^n$ but condition number $1$ in the Euclidean norm: it scales every direction equally. The matrix $\operatorname{diag}(M,1/M)$ has determinant $1$ and condition number $M^2$ for $M\ge1$. Conditioning compares the greatest stretch with the least, rather than their product. An exact theorem about invertibility is not an accuracy guarantee for floating-point computation.
\end{reminder}

\photograph{astronaut}{NASA astronaut Eileen Collins in an orange flight suit beside a helmet.}{Astronaut Eileen Collins. NASA, public domain; via scikit-image. Coordinates and fitted models help connect observations to geometry. A photograph supplies context, not evidence for a numerical model or an endorsement by the person shown.}


\section{Exercises}

\begin{exercise}[$\star$]\label{ex:fibmat}
Let $M=\begin{pmatrix}1&1\\1&0\end{pmatrix}$ and $(F_n)$ the Fibonacci numbers.
\begin{enumerate}[label=(\alph*)]
\item Show that $M^n=\begin{pmatrix}F_{n+1}&F_n\\F_n&F_{n-1}\end{pmatrix}$ for $n\ge1$, and deduce Cassini's identity $F_{n+1}F_{n-1}-F_n^2=(-1)^n$.
\item Diagonalise $M$ and deduce Binet's formula $F_n=\frac{\varphi^n-\psi^n}{\sqrt5}$, where $\varphi=\frac{1+\sqrt5}2$ and $\psi=\frac{1-\sqrt5}2$.
\end{enumerate}
\end{exercise}

\begin{exercise}[$\star$]\label{ex:nilp}
Let $A\in M_n(\C)$ be nilpotent. Show that $\operatorname{tr}A=0$ and $\det(I_n+A)=1$.
\end{exercise}

\begin{exercise}[$\star$]\label{ex:mindet}
Compute the determinant of the $n\times n$ matrix $\left(\min(i,j)\right)_{1\le i,j\le n}$.
\end{exercise}

\begin{exercise}[$\star\star$]\label{ex:annul}
Let $A\in M_n(\R)$ with $A^2-3A+2I_n=0$. Show that $A$ is invertible, compute $A^{-1}$, and show that $A$ is diagonalisable. What are the possible eigenvalues?
\end{exercise}

\begin{exercise}[$\star\star$]\label{ex:abba}
Let $A,B\in M_n(\C)$. Show that $AB$ and $BA$ have the same characteristic polynomial.
\end{exercise}

\begin{exercise}[$\star\star$]\label{ex:rayleigh}
Let $S$ be real symmetric with eigenvalues $\lambda_1\le\dots\le\lambda_n$.
\begin{enumerate}[label=(\alph*)]
\item Show that $\max_{\|x\|=1}x^{\mathsf T}Sx=\lambda_n$ and $\min_{\|x\|=1}x^{\mathsf T}Sx=\lambda_1$.
\item Show that if $x^{\mathsf T}Sx\ge0$ for all $x$, there is a symmetric $R$ with $x^{\mathsf T}Rx\ge0$ for all $x$ and $R^2=S$.
\end{enumerate}
\end{exercise}

\begin{exercise}[$\star\star\star$]\label{ex:hilbert}
\emph{The Hilbert matrix.} Let $H_n=\left(\frac1{i+j-1}\right)_{1\le i,j\le n}$. Show that $H_n$ is the Gram matrix of $1,x,\dots,x^{n-1}$ for the inner product
$\langle f,g\rangle=\int_0^1fg$, and deduce that $\det H_n>0$.
\end{exercise}

\begin{exercise}[$\star\star\star$]\label{ex:codiag}
Let $A,B\in M_n(\C)$ be diagonalisable with $AB=BA$. Show that there is a single invertible $P$ such that $P^{-1}AP$ and $P^{-1}BP$ are both diagonal.
\end{exercise}

\subsection*{Further exercises}

\begin{exercise}[$\star$]\label{ex:twobytwo}
\emph{Complex numbers as matrices.} To $z=a+ib$ associate $M(z)=\begin{pmatrix}a&-b\\b&a\end{pmatrix}$.
\begin{enumerate}[label=(\alph*)]
\item Show that $M(z+w)=M(z)+M(w)$, $M(zw)=M(z)M(w)$ and $\det M(z)=|z|^2$. Deduce $|zw|=|z||w|$ from Theorem~\ref{thm:det}(b).
\item Show that $M(e^{i\theta})$ is the rotation matrix $R_\theta$, and that the identity $R_\theta R_\varphi=R_{\theta+\varphi}$ is the identity $e^{i\theta}e^{i\varphi}=e^{i(\theta+\varphi)}$ of Chapter~2 in disguise: each is equivalent to the addition formulas for $\cos$ and $\sin$.
\item What are the eigenvalues of $M(z)$ over $\C$? Is $M(z)$ diagonalisable over $\R$?
\end{enumerate}
\end{exercise}

\begin{exercise}[$\star\star$]\label{ex:tracerank}
\emph{Trace and rank.}
\begin{enumerate}[label=(\alph*)]
\item Show that $\operatorname{tr}(AB)=\operatorname{tr}(BA)$ for $A\in M_{n,m}$, $B\in M_{m,n}$, and deduce that similar matrices have the same trace. Show that $AB-BA=I_n$ is impossible for $A,B\in M_n(\mathbb K)$ when $\mathbb K=\R$ or $\C$.
\item Show that the rank of the columns of a matrix equals the rank of its rows: $\operatorname{rk}A=\operatorname{rk}A^{\mathsf T}$. (Write $A=BC$ with $B$ having $r=\operatorname{rk}A$ columns.)
\end{enumerate}
\end{exercise}

\begin{exercise}[$\star\star$]\label{ex:hadamard}
\emph{Hadamard's inequality.} Let $A\in M_n(\R)$ have columns $v_1,\dots,v_n$.
\begin{enumerate}[label=(\alph*)]
\item Show that $|\det A|\le\|v_1\|\cdots\|v_n\|$, with equality iff the columns are pairwise orthogonal or one of them is $0$.
\item Deduce that if all entries of $A$ satisfy $|a_{ij}|\le1$, then $|\det A|\le n^{n/2}$. Find a $2\times2$ and a $4\times4$ matrix with entries $\pm1$ attaining the bound.
\end{enumerate}
\end{exercise}

\begin{exercise}[$\star\star$]\label{ex:sylvester}
\emph{Sylvester's determinant identity.} For $A\in M_{m,n}(\mathbb K)$ and $B\in M_{n,m}(\mathbb K)$, show that $\det(I_m+AB)=\det(I_n+BA)$. (Multiply the block matrix $\begin{pmatrix}I_m&-A\\B&I_n\end{pmatrix}$ on the left, then on the right, by $\begin{pmatrix}I_m&A\\0&I_n\end{pmatrix}$.) Deduce that the non-zero eigenvalues of $AB$ and $BA$ coincide, and compute $\det(I_n+uv^{\mathsf T})$ for column vectors $u,v$.
\end{exercise}

\begin{exercise}[$\star$]\label{ex:fibch}
\emph{Fibonacci by Cayley--Hamilton.} With $M$ as in Exercise~\ref{ex:fibmat}, use $M^2=M+I_2$ to show that $M^n=F_nM+F_{n-1}I_2$ for $n\ge1$, and deduce $F_{m+n}=F_mF_{n+1}+F_{m-1}F_n$ and $F_{2n}=F_n(F_{n+1}+F_{n-1})$. Which Fibonacci numbers divide $F_{2n}$? Which divide $F_{kn}$?
\end{exercise}

\begin{exercise}[$\star\star$]\label{ex:intinverse}
\emph{Integer matrices.} Let $A\in M_n(\Z)$.
\begin{enumerate}[label=(\alph*)]
\item Show that $A$ has an inverse with integer entries if and only if $\det A=\pm1$.
\item If $d=\det A\ne0$, show that $d\,\Z^n\subseteq A\Z^n$: every vector of $\Z^n$ divisible by $d$ is an integer combination of the columns of $A$.
\item Two bases of $\Z^2$ (pairs $u,v$ with $\Z u+\Z v=\Z^2$) span parallelograms of the same area. Which?
\end{enumerate}
\end{exercise}

\begin{exercise}[$\star\star\star$]\label{ex:stochastic}
\emph{A random surfer.} A matrix $G\in M_n(\R)$ is \emph{positive stochastic} if all its entries are $>0$ and each column sums to $1$. (The Google matrix of Figure~\ref{fig:pagerank} is one: $g_{ij}$ is the probability that a surfer on page $j$ moves to page $i$.)
\begin{enumerate}[label=(\alph*)]
\item Show that $1$ is an eigenvalue of $G$, and that every complex eigenvalue satisfies $|\lambda|\le1$.
\item Let $c=\min_{i,j}g_{ij}$ and let $\|x\|_1=\sum|x_i|$. Show that if $\sum_ix_i=0$ then $\|Gx\|_1\le(1-nc)\|x\|_1$.
\item Deduce that there is exactly one probability vector $\pi$ (non-negative entries summing to $1$) with $G\pi=\pi$, that all its entries are positive, and that for every probability vector $x_0$ the sequence $G^kx_0$ converges to $\pi$, with $\|G^kx_0-\pi\|_1\le2(1-nc)^k$.
\item Compute $\pi$ for the four-page web of Figure~\ref{fig:pagerank}.
\end{enumerate}
\end{exercise}

\begin{figure}[ht]
\centering
\fig{ch07-pagerank}
\caption{Left: a web of four pages and its links. Page $1$ links to $2$ and $3$; page $2$ to $3$; page $3$ to $1$; page $4$ to $1$ and $3$. Middle: the matrix $P$ of the surfer who clicks a random link ($p_{ij}=1/\text{out-degree}(j)$ if $j$ links to $i$), and the Google matrix $G=0.85P+0.15\cdot\frac14J$, where $J$ is the all-ones matrix: with probability $0.15$ the surfer jumps to a random page instead. Right: the entries of $G^kx_0$ for $x_0=(\frac14,\frac14,\frac14,\frac14)$; they converge geometrically to the PageRank vector $\pi$, with page $3$ ranked first and page $4$ last. Page~$4$ has no incoming link: its only visitors arrive by the random jump.}
\label{fig:pagerank}
\end{figure}

\begin{exercise}[$\star\star\star$]\label{ex:cauchydet}
\emph{Cauchy's determinant and the Hilbert matrix.} Let $x_1,\dots,x_n,y_1,\dots,y_n$ be scalars with all $x_i+y_j\ne0$.
\begin{enumerate}[label=(\alph*)]
\item Show that
\[
\det\left(\frac1{x_i+y_j}\right)_{1\le i,j\le n}=\frac{\prod_{1\le i<j\le n}(x_j-x_i)(y_j-y_i)}{\prod_{1\le i,j\le n}(x_i+y_j)}.
\]
\item Deduce $\det H_n=\dfrac{c_n^4}{c_{2n}}$ with $c_n=\prod_{k=1}^{n-1}k!$, and check $\det H_2=\frac1{12}$, $\det H_3=\frac1{2160}$.
\item Show that $\det H_n<4^{-(n-1)^2}$ for $n\ge2$: the determinant decays faster than any geometric sequence.
\end{enumerate}
\end{exercise}

\begin{challenge}
\textbf{Three doors in the wall.} (1) \emph{Perron--Frobenius.} Exercise~\ref{ex:stochastic} assumed all entries positive. Show that for any matrix with positive entries (not necessarily stochastic) there is an eigenvalue $r>0$ which is a simple root of $\chi$, has an eigenvector with positive entries, and satisfies $|\lambda|<r$ for every other eigenvalue. Then ask what survives when some entries are $0$: the Google matrix adds $0.15J/n$ precisely to avoid this question.
(2) \emph{Hadamard's bound is sharp when it is.} A matrix with entries $\pm1$ and $|\det|=n^{n/2}$ is a \emph{Hadamard matrix}; show that it exists only for $n=1,2$ or $n\equiv0\pmod 4$. Whether it exists for \emph{every} multiple of $4$ is open; as of October 2026 the smallest undecided order is $2396$ (see the notes and sources).
(3) \emph{Every matrix is a product of two symmetric matrices} (Frobenius, 1910), over any field. Try $n=2$ by hand, then find why the general case is hard. Volume~2 gives the tools.
\end{challenge}

\section{Hints}

\begin{hint}{ex:fibmat}
(a) Induction, then take determinants. (b) The eigenvalues of $M$ are the roots of $X^2-X-1$.
\end{hint}
\begin{hint}{ex:nilp}
If $Av=\lambda v$ with $v\ne0$, then $A^kv=\lambda^kv$.
\end{hint}
\begin{hint}{ex:mindet}
Subtract each row from the next one, starting from the bottom.
\end{hint}
\begin{hint}{ex:annul}
$A(3I_n-A)=2I_n$, and $X^2-3X+2=(X-1)(X-2)$.
\end{hint}
\begin{hint}{ex:abba}
If $A$ is invertible, $BA=A^{-1}(AB)A$. Otherwise, replace $A$ by $A+tI_n$, and view both characteristic polynomials as polynomials in $t$.
\end{hint}
\begin{hint}{ex:rayleigh}
(a) Write $x=Py$ in an orthonormal eigenbasis: $x^{\mathsf T}Sx=\sum\lambda_iy_i^2$. (b) Take square roots of the eigenvalues.
\end{hint}
\begin{hint}{ex:hilbert}
$\frac1{i+j-1}=\int_0^1x^{i-1}x^{j-1}\dd x$. For $c\ne0$, what is $c^{\mathsf T}H_nc$?
\end{hint}
\begin{hint}{ex:codiag}
$B$ maps each eigenspace of $A$ into itself. Diagonalise $B$ inside each one.
\end{hint}
\begin{hint}{ex:twobytwo}
(a) Multiply out. (b) $e^{i\theta}=\cos\theta+i\sin\theta$. (c) $\chi=X^2-2aX+a^2+b^2$.
\end{hint}
\begin{hint}{ex:tracerank}
(a) Both traces are $\sum_{i,j}a_{ij}b_{ji}$. (b) If $b_1,\dots,b_r$ is a basis of the column space, each column of $A$ is a combination of them; write that as $A=BC$ and transpose.
\end{hint}
\begin{hint}{ex:hadamard}
(a) In Theorem~\ref{thm:gram}, $d_k\le\|v_k\|$. (b) Each column has norm $\le\sqrt n$. For $n=2$ take $H=\begin{pmatrix}1&1\\1&-1\end{pmatrix}$, for $n=4$ the block matrix $\begin{pmatrix}H&H\\H&-H\end{pmatrix}$.
\end{hint}
\begin{hint}{ex:sylvester}
The two products are block triangular, with blocks $I_m+AB$ and $I_n+BA$ on their diagonals. For the eigenvalues, replace $A$ by $-A/\lambda$.
\end{hint}
\begin{hint}{ex:fibch}
Induction for $M^n$; then compare the entries of $M^{m+n}=M^mM^n$ using Exercise~\ref{ex:fibmat}(a).
\end{hint}
\begin{hint}{ex:intinverse}
(a) $\det A\det A^{-1}=1$ with both factors integers; conversely use the adjugate. (b) $A\operatorname{adj}A=dI$. (c) The change-of-basis matrix is in $M_2(\Z)$ with integer inverse.
\end{hint}
\begin{hint}{ex:stochastic}
(a) The row vector $\mathbf 1=(1,\dots,1)$ satisfies $\mathbf 1G=\mathbf 1$; use $\det(G^{\mathsf T}-\lambda I)=\det(G-\lambda I)$. For $|\lambda|\le1$ consider a left eigenvector $y$ and the index where $|y_i|$ is largest. (b) Write each column as $c\mathbf 1^{\mathsf T}+h_j$ with $h_j\ge0$. (c) The sequence $G^kx_0$ is Cauchy for $\|\cdot\|_1$; two fixed probability vectors have a difference of sum $0$.
\end{hint}
\begin{hint}{ex:cauchydet}
(a) Subtract the last column from each other column, factor $y_n-y_j$ out of column $j$ and $\frac1{x_i+y_n}$ out of row $i$; then do the same with rows, and expand along the last column, which has become $(0,\dots,0,1)^{\mathsf T}$. (b) $x_i=i$, $y_j=j-1$. (c) $c_{2n}/c_n^4=\prod_{k=1}^{n-1}\binom{2k}k^2(2k+1)$, and $\binom{2k}k\ge\frac{4^k}{2\sqrt k}$ (induction on $k$).
\end{hint}

\section{Solutions}

\begin{solution}{ex:fibmat}
\idea A two-term recurrence is a matrix power, so identities about $F_{n\pm1},F_n$ are identities about $M^n$; Cassini's left side looks like a $2\times2$ determinant, and $\det M^n=(\det M)^n$ is free. For Binet, diagonalising $M$ shows every entry of $M^n$ is a combination of $\varphi^n,\psi^n$; two initial values fix the coefficients.
\writeup (a) True for $n=1$ (with $F_0=0$), and $M^{n+1}=M^nM=\begin{pmatrix}F_{n+1}+F_n&F_{n+1}\\F_n+F_{n-1}&F_n\end{pmatrix}$ has the right form.
Taking determinants (Theorem~\ref{thm:det}(b)), $F_{n+1}F_{n-1}-F_n^2=(\det M)^n=(-1)^n$.

(b) $\chi_M=X^2-X-1$ has roots $\varphi\ne\psi$, so $M$ is diagonalisable (Theorem~\ref{thm:diag}): $M=PDP^{-1}$ and $M^n=PD^nP^{-1}$. Each entry of $M^n$ is therefore of the form
$\alpha\varphi^n+\beta\psi^n$; for $F_n$, the values $F_0=0$, $F_1=1$ give $\alpha+\beta=0$, $\alpha\varphi+\beta\psi=1$, so $\alpha=-\beta=\frac1{\varphi-\psi}=\frac1{\sqrt5}$.
\end{solution}

\begin{solution}{ex:nilp}
\idea Trace and determinant are read off the eigenvalues (Proposition~\ref{prop:charpoly}), so find the eigenvalues of a nilpotent matrix: apply $A^k=0$ to an eigenvector. Then $I_n+A$ has its eigenvalues shifted by $1$.
\writeup If $A^k=0$ and $Av=\lambda v$ with $v\ne0$, then $0=A^kv=\lambda^kv$, so $\lambda=0$. Over $\C$, $\chi_A$ splits, and all its roots are $0$: $\chi_A=X^n$. Hence
$\operatorname{tr}A=0$, and $\chi_{I+A}(X)=\det((X-1)I-A)=\chi_A(X-1)=(X-1)^n$, whose constant term $(-1)^n$ equals $(-1)^n\det(I_n+A)$: so $\det(I_n+A)=1$.
\end{solution}

\begin{solution}{ex:mindet}
\idea The entries $\min(i,j)$ grow by $1$ from one row to the next, but only in the columns $j\ge i$. So the difference of consecutive rows is a row of $0$s then $1$s, and row operations (Theorem~\ref{thm:det}(d)) should make the matrix triangular. Work from the bottom up so that each subtraction uses an original row.
\writeup For $i=n,n-1,\dots,2$ in that order, subtract row $i-1$ from row $i$. The new row $i$ has entries $\min(i,j)-\min(i-1,j)$, which is $1$ for $j\ge i$ and $0$ for $j<i$.
The matrix becomes upper triangular with ones on the diagonal: the determinant is $1$.
\end{solution}

\begin{solution}{ex:annul}
\idea An annihilating polynomial with non-zero constant term gives the inverse: isolate the constant and factor $A$ out of the rest. The polynomial factors as $(X-1)(X-2)$, split with simple roots, which is exactly the criterion of Theorem~\ref{thm:diag}.
\writeup $A(3I_n-A)=2I_n$, so $A$ is invertible with $A^{-1}=\frac12(3I_n-A)$. The polynomial $X^2-3X+2=(X-1)(X-2)$ is split with simple roots and cancels $A$,
so $A$ is diagonalisable, with eigenvalues in $\{1,2\}$ (if $Av=\lambda v$, $v\ne0$, then $0=(A^2-3A+2I)v=(\lambda^2-3\lambda+2)v$). Each value can occur: $I_n$ has only $1$, and $2I_n$ only $2$. For $n\ge2$, a diagonal matrix with both $1$ and $2$ has both eigenvalues.
\end{solution}

\begin{solution}{ex:abba}
\idea When $A$ is invertible, $BA$ is a conjugate of $AB$, and similar matrices have the same $\chi$. The general case is a \emph{density argument}: perturb $A$ to $A+tI$, invertible for all but finitely many $t$, and use that a polynomial identity in $t$ true at infinitely many points is true everywhere.
\writeup If $A$ is invertible, $BA=A^{-1}(AB)A$ is similar to $AB$, so they have the same characteristic polynomial (Proposition~\ref{prop:charpoly}). In general, $A+tI_n$ is invertible for all but finitely many $t\in\C$
(those $-t$ that are eigenvalues of $A$). For such $t$ and every $x$,
\[
\begin{aligned}&\det\left(xI_n-(A+tI_n)B\right)\eqbr=\det\left(xI_n-B(A+tI_n)\right).\end{aligned}
\]
For fixed $x$, both sides are polynomials in $t$ that agree at infinitely many points, so they are equal (Proposition~\ref{prop:rootsfactor}), in particular at $t=0$.
\end{solution}

\begin{solution}{ex:rayleigh}
\idea Questions about a symmetric matrix become questions about a diagonal one in an orthonormal eigenbasis (Theorem~\ref{thm:spectral}), and orthogonal changes of basis preserve unit vectors. Then $x^{\mathsf T}Sx$ is a weighted average $\sum\lambda_iy_i^2$ with weights $y_i^2$ summing to $1$. For (b), take square roots of the eigenvalues, which is possible since they are $\ge0$ by (a).
\writeup (a) Write $S=PDP^{\mathsf T}$ and $y=P^{\mathsf T}x$; $P$ is orthogonal so $\|y\|=\|x\|=1$, and $x^{\mathsf T}Sx=y^{\mathsf T}Dy=\sum\lambda_iy_i^2$. This lies between $\lambda_1\sum y_i^2=\lambda_1$
and $\lambda_n$, and the bounds are attained at unit eigenvectors.

(b) By (a) all $\lambda_i\ge0$. Let $R=P\sqrt DP^{\mathsf T}$ where $\sqrt D=\operatorname{diag}(\sqrt{\lambda_i})$. Then $R$ is symmetric, $R^2=PDP^{\mathsf T}=S$, and its eigenvalues
$\sqrt{\lambda_i}\ge0$ give $x^{\mathsf T}Rx\ge0$ by (a).
\end{solution}

\begin{solution}{ex:hilbert}
\idea The number $\frac1{i+j-1}$ is the integral $\int_0^1x^{i-1}x^{j-1}\dd x$, and seeing it that way turns the matrix into a Gram matrix. Gram matrices of independent vectors are positive definite, because $c^{\mathsf T}Gc$ is a squared norm; positive definite symmetric matrices have positive eigenvalues, hence positive determinant.
\writeup For the scalar product $\langle f,g\rangle=\int_0^1fg$,
\[
\langle x^{i-1},x^{j-1}\rangle=\int_0^1x^{i+j-2}\dd x=\frac1{i+j-1}.
\]
For a Gram matrix $G$ with entries $\langle e_i,e_j\rangle$, of linearly independent vectors and $c\in\R^n\setminus\{0\}$,
$c^{\mathsf T}Gc=\left\|\sum c_ie_i\right\|^2>0$, because $\sum c_ie_i\ne0$. So $G$ is symmetric with all eigenvalues positive (by Exercise~\ref{ex:rayleigh}(a), the least one is $\min_{\|c\|=1}c^{\mathsf T}Gc>0$), and $\det G$, their product, is positive.
The monomials $1,x,\dots,x^{n-1}$ are linearly independent (a non-zero polynomial has finitely many roots), so $\det H_n>0$.
\end{solution}

\begin{solution}{ex:codiag}
\idea Commuting matrices preserve each other's eigenspaces. So decompose the space into eigenspaces of $A$ and diagonalise $B$ inside each one. That $B$ stays diagonalisable on a stable subspace follows from the polynomial criterion: a polynomial that kills $B$ kills its restriction.
\writeup Let $\lambda_1,\dots,\lambda_r$ be the distinct eigenvalues of $A$ and $E_k=\ker(A-\lambda_kI_n)$; since $A$ is diagonalisable, $\C^n=E_1\oplus\dots\oplus E_r$ (Theorem~\ref{thm:diag}).
Each $E_k$ is stable under $B$: if $Av=\lambda_kv$ then $A(Bv)=BAv=\lambda_kBv$. A polynomial with simple roots that cancels $B$ also cancels its restriction $B_k$ to $E_k$,
so $B_k$ is diagonalisable: $E_k$ has a basis of eigenvectors of $B$, which are automatically eigenvectors of $A$ (for $\lambda_k$). Gathering these bases gives a basis of $\C^n$ made of common eigenvectors; let $P$ have them as columns.
\end{solution}


\begin{solution}{ex:twobytwo}
\idea Everything is a direct computation; the point is to notice that the computation of $M(z)M(w)$ \emph{is} the computation of $(a+ib)(c+id)$, with the same four products in the same places. So the complex numbers are a copy of a family of $2\times2$ matrices, and matrix facts (determinants multiply) become facts about complex numbers (moduli multiply).
\writeup (a) Additivity is clear. For $z=a+ib$, $w=c+id$,
\[
M(z)M(w)=\begin{pmatrix}ac-bd&-(ad+bc)\\ad+bc&ac-bd\end{pmatrix}=M\big((ac-bd)+i(ad+bc)\big)=M(zw),
\]
and $\det M(z)=a^2+b^2=|z|^2$. Hence $|zw|^2=\det M(zw)=\det M(z)\det M(w)=|z|^2|w|^2$.

(b) Since $e^{i\theta}=\cos\theta+i\sin\theta$,
\[
M(e^{i\theta})=\begin{pmatrix}\cos\theta&-\sin\theta\\\sin\theta&\cos\theta\end{pmatrix}=R_\theta.
\] Since $M$ is injective and multiplicative, $R_\theta R_\varphi=R_{\theta+\varphi}$ holds iff $M(e^{i\theta}e^{i\varphi})=M(e^{i(\theta+\varphi)})$ iff $e^{i\theta}e^{i\varphi}=e^{i(\theta+\varphi)}$. Reading the first column of either identity gives $\cos(\theta+\varphi)=\cos\theta\cos\varphi-\sin\theta\sin\varphi$ and $\sin(\theta+\varphi)=\sin\theta\cos\varphi+\cos\theta\sin\varphi$, and conversely these two formulas give the identity.

(c) $\chi_{M(z)}=(X-a)^2+b^2$, with roots $a\pm ib=z,\bar z$; the eigenvectors are $(1,-i)$ for $z$ and $(1,i)$ for $\bar z$ (check: $M(z)\binom1{-i}=\binom{a+ib}{b-ia}=z\binom1{-i}$). Over $\R$, if $b\ne0$ there is no real eigenvalue at all, so $M(z)$ is not diagonalisable; if $b=0$ it is the scalar matrix $aI_2$. A rotation of the plane by an angle that is not a multiple of $\pi$ fixes no direction, which is the geometric content.
\end{solution}

\begin{solution}{ex:tracerank}
\idea (a) Write the trace as a double sum and swap the order. (b) The rank is the number of columns one needs to generate all of them; writing ``each column is a combination of $r$ chosen columns'' as a matrix product $A=BC$ and transposing turns a statement about columns into one about rows.
\writeup (a) $\operatorname{tr}(AB)=\sum_i\sum_ja_{ij}b_{ji}=\sum_j\sum_ib_{ji}a_{ij}=\operatorname{tr}(BA)$. Then $\operatorname{tr}(P^{-1}AP)=\operatorname{tr}(APP^{-1})=\operatorname{tr}A$, so the trace of a linear map is defined. If $AB-BA=I_n$, taking traces gives $0=n$, impossible in $\R$ or $\C$. (In infinite dimension the identity is possible: on polynomials, $(\frac{d}{dx}\circ x-x\circ\frac{d}{dx})P=P$. This is Heisenberg's commutation relation, and the trace argument is why quantum mechanics needs infinitely many dimensions.)

(b) Let $r=\operatorname{rk}A$ be the dimension of the column space and $b_1,\dots,b_r$ a basis of it; let $B\in M_{n,r}$ have these as columns. Each column $a_j$ of $A$ is $\sum_kc_{kj}b_k$, which says $A=BC$ with $C=(c_{kj})\in M_{r,m}$. Transposing, $A^{\mathsf T}=C^{\mathsf T}B^{\mathsf T}$, whose columns are combinations of the $r$ columns of $C^{\mathsf T}$: so $\operatorname{rk}A^{\mathsf T}\le r=\operatorname{rk}A$. Applying this to $A^{\mathsf T}$ gives the reverse inequality.
\end{solution}

\begin{solution}{ex:hadamard}
\idea Theorem~\ref{thm:gram} writes $|\det A|$ as a product of heights $d_k$, and a height is at most the length of the edge: that is the whole proof, and equality means every edge is its own height, that is, perpendicular to the previous ones. For (b), a vector with $n$ entries of size $\le1$ has length $\le\sqrt n$.
\writeup (a) If the columns are dependent, $\det A=0\le\prod\|v_k\|$, with equality iff some $v_k=0$ (if all $v_k\ne0$ the product is positive). Otherwise, by Theorem~\ref{thm:gram}, $|\det A|=\prod_kd_k$ with $d_k=\|v_k-p_k(v_k)\|$, $p_k$ the projection on $F_k=\operatorname{Span}(v_1,\dots,v_{k-1})$. Pythagoras gives $\|v_k\|^2=\|p_k(v_k)\|^2+d_k^2$, so $d_k\le\|v_k\|$, and equality in the product (all factors positive) forces $p_k(v_k)=0$ for every $k$, that is $v_k\perp F_k$: the columns are pairwise orthogonal. Conversely, pairwise orthogonal non-zero columns have $p_k(v_k)=0$ and give equality.

(b) $\|v_k\|^2=\sum_ia_{ik}^2\le n$, so $|\det A|\le n^{n/2}$. Equality needs entries $\pm1$ and orthogonal columns. $H=\begin{pmatrix}1&1\\1&-1\end{pmatrix}$ has $|\det H|=2=2^{2/2}$, and $H_4=\begin{pmatrix}H&H\\H&-H\end{pmatrix}$ has orthogonal columns of norm $2$, so $|\det H_4|=2^4=16=4^{4/2}$. (The construction $H\mapsto\begin{pmatrix}H&H\\H&-H\end{pmatrix}$ doubles the order, so Hadamard matrices exist for every power of $2$; this is Sylvester's construction of 1867.)
\end{solution}

\begin{solution}{ex:sylvester}
\idea A block matrix with an identity block can be made block triangular by block row or column operations, and two different ways of doing it expose $I_m+AB$ and $I_n+BA$ as the only non-trivial diagonal blocks. The eigenvalue statement is the identity with $A$ scaled by $-1/\lambda$, and $\det(I+uv^{\mathsf T})$ is the case $m=n$, $n=1$ in reverse: a $1\times1$ determinant.
\writeup Let $M=\begin{pmatrix}I_m&-A\\B&I_n\end{pmatrix}$ and $T=\begin{pmatrix}I_m&A\\0&I_n\end{pmatrix}$, with $\det T=1$ (triangular). Block multiplication gives
\[
TM=\begin{pmatrix}I_m+AB&0\\B&I_n\end{pmatrix},\qquad MT=\begin{pmatrix}I_m&0\\B&I_n+BA\end{pmatrix}.
\]
For a block lower triangular matrix $\begin{pmatrix}X&0\\Y&Z\end{pmatrix}$ the determinant is $\det X\det Z$: indeed the matrix equals
\[
\begin{pmatrix}X&0\\0&I\end{pmatrix}\begin{pmatrix}I&0\\Y&I\end{pmatrix}\begin{pmatrix}I&0\\0&Z\end{pmatrix},
\]
the middle factor is triangular with determinant $1$, and $\det\begin{pmatrix}X&0\\0&I\end{pmatrix}=\det X$ because, as a function of the columns of $X$, it is multilinear, alternating and equal to $1$ at $X=I$ (Theorem~\ref{thm:det}(a)); likewise for $Z$. So $\det M=\det(I_m+AB)=\det(I_n+BA)$.

For $\lambda\ne0$, applying this to $-A/\lambda$ and $B$:
\[
\det(\lambda I_m-AB)=\lambda^m\det\left(I_m-\tfrac1\lambda AB\right)=\lambda^m\det\left(I_n-\tfrac1\lambda BA\right)=\lambda^{m-n}\det(\lambda I_n-BA),
\]
so $\lambda^n\chi_{AB}(\lambda)=\lambda^m\chi_{BA}(\lambda)$ for all $\lambda\ne0$, hence as polynomials: $\chi_{AB}$ and $\chi_{BA}$ differ by a power of $X$, and their non-zero roots coincide with multiplicities. (For $m=n$ this is Exercise~\ref{ex:abba} again.) Finally, with $A=u$ and $B=v^{\mathsf T}$ ($n=1$ on the small side), $\det(I_n+uv^{\mathsf T})=\det(1+v^{\mathsf T}u)=1+v^{\mathsf T}u=1+\sum_i u_iv_i$.
\end{solution}

\begin{solution}{ex:fibch}
\idea Cayley--Hamilton reduces every power of $M$ to the plane spanned by $I$ and $M$, and the coefficients are forced to be Fibonacci numbers by the recurrence. Then $M^{m+n}=M^mM^n$ is a product of two such expressions, and comparing coefficients gives the addition formula for free. Everything in the exercise is this one idea.
\writeup For $n=1$, $M=F_1M+F_0I$. If $M^n=F_nM+F_{n-1}I$, then $M^{n+1}=F_nM^2+F_{n-1}M=F_n(M+I)+F_{n-1}M=F_{n+1}M+F_nI$. Now
\[
\begin{aligned}
M^{m+n}&=(F_mM+F_{m-1}I)(F_nM+F_{n-1}I)\\
&=F_mF_n(M+I)+(F_mF_{n-1}+F_{m-1}F_n)M+F_{m-1}F_{n-1}I\\
&=(F_mF_{n+1}+F_{m-1}F_n)M+(F_mF_n+F_{m-1}F_{n-1})I ,
\end{aligned}
\]
and also $M^{m+n}=F_{m+n}M+F_{m+n-1}I$. Since $M$ is not a scalar matrix, $I$ and $M$ are linearly independent and the coefficients match: $F_{m+n}=F_mF_{n+1}+F_{m-1}F_n$. With $m=n$, $F_{2n}=F_n(F_{n+1}+F_{n-1})$; the second factor is the Lucas number $L_n$. So $F_n$ divides $F_{2n}$, and by induction on $k$, $F_n$ divides $F_{kn}$: $F_{(k+1)n}=F_{kn}F_{n+1}+F_{kn-1}F_n$ with both terms divisible by $F_n$. (Exercise~\ref{ex:fibgcd} in Chapter~1 proves the stronger identity $\gcd(F_m,F_n)=F_{\gcd(m,n)}$.)
\end{solution}

\begin{solution}{ex:intinverse}
\idea The adjugate formula shows that $A^{-1}$ has entries in $\frac1{\det A}\Z$, so the only obstacle to an integer inverse is the determinant; and $\det A\cdot\det A^{-1}=1$ with two integers leaves only $\pm1$. For (c), a change of basis of $\Z^2$ is an integer matrix with integer inverse, so its determinant is $\pm1$ and it preserves areas.
\writeup (a) If $A^{-1}\in M_n(\Z)$, then $\det A\cdot\det A^{-1}=1$ with both factors in $\Z$, so $\det A=\pm1$. Conversely, if $\det A=\pm1$, then $A^{-1}=\pm\operatorname{adj}A$ (Proposition~\ref{prop:adjugate}), and the cofactors are determinants of integer matrices, hence integers.

(b) For $w\in\Z^n$, $dw=A\operatorname{adj}(A)w$ and $\operatorname{adj}(A)w\in\Z^n$: so $dw$ is an integer combination of the columns of $A$.

(c) If $(u,v)$ and $(u',v')$ are bases of $\Z^2$, write $u'=au+cv$, $v'=bu+dv$ with $P=\begin{pmatrix}a&b\\c&d\end{pmatrix}\in M_2(\Z)$; expressing $u,v$ in terms of $u',v'$ gives an integer matrix which is $P^{-1}$, so $\det P=\pm1$ by (a), and $|\det(u',v')|=|\det P|\,|\det(u,v)|=|\det(u,v)|$. Comparing with the basis $(e_1,e_2)$: every basis of $\Z^2$ spans a parallelogram of area exactly $1$. Conversely, if $u,v\in\Z^2$ span area $1$, the matrix $(u\ v)$ has determinant $\pm1$, so by (a) its inverse is integral, and $e_1,e_2$ are integer combinations of $u,v$: it is a basis. (Pick's theorem, in Volume~1, says the same thing in pictures: such a parallelogram contains no lattice point other than its corners.)
\end{solution}

\begin{solution}{ex:stochastic}
\idea Column sums equal to $1$ say that the row vector of ones is a left eigenvector for $1$. The contraction in (b) is the heart: subtracting from each column the constant vector $c\mathbf 1$ leaves non-negative columns of sum $1-nc$, and the constant part is invisible to vectors of sum $0$. After that, (c) is the standard route from ``contraction'' to ``unique fixed point, reached geometrically'': Cauchy sequences converge, and two fixed points would have to be at distance $0$.
\writeup (a) $\mathbf 1^{\mathsf T}G=\mathbf 1^{\mathsf T}$ because each column of $G$ sums to $1$; so $G^{\mathsf T}\mathbf 1=\mathbf 1$, $1$ is an eigenvalue of $G^{\mathsf T}$, and $\chi_{G^{\mathsf T}}=\chi_G$ (Theorem~\ref{thm:det}(b)), so of $G$. If $\lambda\in\C$ is an eigenvalue of $G$, it is one of $G^{\mathsf T}$: there is $y\ne0$ with $y^{\mathsf T}G=\lambda y^{\mathsf T}$, that is $\lambda y_j=\sum_iy_ig_{ij}$ for every $j$. Choose $j$ with $|y_j|$ maximal; then $|\lambda||y_j|\le\sum_i|y_i|g_{ij}\le|y_j|\sum_ig_{ij}=|y_j|$, so $|\lambda|\le1$.

(b) Write column $j$ of $G$ as $g_j=c\mathbf 1+h_j$ with $h_j\ge0$ entrywise and $\sum_ih_{ij}=1-nc$. For $x$ with $\sum_jx_j=0$,
\[
Gx=\sum_jx_jg_j=c\Big(\sum_jx_j\Big)\mathbf 1+\sum_jx_jh_j=\sum_jx_jh_j ,
\]
so $\|Gx\|_1\le\sum_j|x_j|\,\|h_j\|_1=(1-nc)\|x\|_1$. (Note $nc\le1$, since a column has $n$ entries $\ge c$ summing to $1$.) Put $q=1-nc<1$.

(c) $G$ maps probability vectors to probability vectors (non-negative entries, and $\mathbf 1^{\mathsf T}Gx=\mathbf 1^{\mathsf T}x=1$). Let $x_k=G^kx_0$. Then $x_{k+1}-x_k=G^k(x_1-x_0)$ with $x_1-x_0$ of sum $0$, so $\|x_{k+1}-x_k\|_1\le q^k\|x_1-x_0\|_1\le2q^k$ by (b) applied $k$ times (each $G^jx$ keeps sum $0$). Hence $\|x_l-x_k\|_1\le2\sum_{j\ge k}q^j=\frac{2q^k}{1-q}$ for $l>k$: the sequence is Cauchy in each coordinate, so it converges (Chapter~5) to some $\pi$, which is a probability vector (limits of non-negative numbers summing to $1$), and $G\pi=\lim Gx_k=\lim x_{k+1}=\pi$. If $\pi'$ is another fixed probability vector, $\pi-\pi'$ has sum $0$ and $\|\pi-\pi'\|_1=\|G(\pi-\pi')\|_1\le q\|\pi-\pi'\|_1$, forcing $\pi=\pi'$. Each entry is positive: $\pi_i=\sum_jg_{ij}\pi_j\ge c\sum_j\pi_j=c>0$. Finally $x_k-\pi=G^k(x_0-\pi)$ with $x_0-\pi$ of sum $0$ and $\|x_0-\pi\|_1\le2$, so $\|x_k-\pi\|_1\le2q^k$.

(d) Here $c=0.15/4=0.0375$ and $q=0.85$. The last row of $G$ is constant, equal to $0.0375$, so $\pi_4=0.0375\sum_j\pi_j=0.0375$ exactly: page~$4$ gets precisely the random-jump traffic. Solving $(G-I)\pi=0$ with $\sum\pi_j=1$ (three unknowns left) gives
\[
\pi=\left(\tfrac{2687}{7076},\ \tfrac{56293}{283040},\ \tfrac{108653}{283040},\ \tfrac{3}{80}\right)\approx(0.3797,\ 0.1989,\ 0.3839,\ 0.0375):
\]
page~$3$ ranks first, just ahead of page~$1$, which it links to. Figure~\ref{fig:pagerank} shows the power iteration; the actual convergence factor is $0.60$, the modulus of the other eigenvalues ($\chi_G=X(X-1)(X^2+0.85X+0.36125)$), better than the guaranteed $q=0.85$.
\end{solution}

\begin{solution}{ex:cauchydet}
\idea The entries $\frac1{x_i+y_j}$ have a product structure that row and column operations reveal: the difference of two such entries factors as $\frac{(y_n-y_j)}{(x_i+y_j)(x_i+y_n)}$. Clearing one column of $1$s and one row of $0$s costs a known product of factors and leaves the same determinant one size smaller: induction. The Hilbert matrix is the case $x_i=i$, $y_j=j-1$, and then every product is a product of factorials.
\writeup (a) Let $D_n$ be the determinant. Subtract column $n$ from each column $j<n$: entry $(i,j)$ becomes
\[
\frac1{x_i+y_j}-\frac1{x_i+y_n}=\frac{y_n-y_j}{(x_i+y_j)(x_i+y_n)}.
\]
Factor $y_n-y_j$ out of column $j$ ($j<n$) and $\frac1{x_i+y_n}$ out of row $i$ (all $i$): the matrix now has entries $\frac1{x_i+y_j}$ for $j<n$ and $1$ in column $n$. Subtract row $n$ from each row $i<n$: entry $(i,j)$, $j<n$, becomes $\frac{x_n-x_i}{(x_i+y_j)(x_n+y_j)}$, and column $n$ becomes $(0,\dots,0,1)^{\mathsf T}$. Factor $x_n-x_i$ out of row $i<n$ and $\frac1{x_n+y_j}$ out of column $j<n$. Expanding along the last column,
\[
D_n=\frac{\prod_{j<n}(y_n-y_j)\prod_{i<n}(x_n-x_i)}{\prod_{i\le n}(x_i+y_n)\prod_{j<n}(x_n+y_j)}\;D_{n-1},
\]
and $D_1=\frac1{x_1+y_1}$. Multiplying these relations for $n,n-1,\dots,2$ gives the formula: the numerators assemble $\prod_{i<j}(x_j-x_i)(y_j-y_i)$ and the denominators every $x_i+y_j$ exactly once.

(b) With $x_i=i$, $y_j=j-1$, $x_i+y_j=i+j-1$. For fixed $j$, $\prod_{i<j}(j-i)=(j-1)!$, so $\prod_{i<j}(x_j-x_i)=\prod_{j=1}^n(j-1)!=c_n$, and the same for the $y$'s: the numerator is $c_n^2$. The denominator is $\prod_{i=1}^n\prod_{j=1}^n(i+j-1)=\prod_{i=1}^n\frac{(i+n-1)!}{(i-1)!}=\frac{\prod_{m=n}^{2n-1}m!}{c_n}=\frac{c_{2n}}{c_n^2}$. So $\det H_n=c_n^4/c_{2n}$. Check: $c_2=1$, $c_4=1!\,2!\,3!=12$, $\det H_2=\frac1{12}$; $c_3=2$, $c_6=1!\,2!\,3!\,4!\,5!=34560$, $\det H_3=\frac{16}{34560}=\frac1{2160}$.

(c) Grouping $m!$ for $m=2k$ and $2k+1$, $c_{2n}=\prod_{k=1}^{n-1}(2k)!(2k+1)!$ (the factor $1!$ is $1$), so
\[
\frac{c_{2n}}{c_n^4}=\prod_{k=1}^{n-1}\frac{(2k)!\,(2k+1)!}{(k!)^4}=\prod_{k=1}^{n-1}\binom{2k}k^2(2k+1).
\]
Now $\binom{2k}k\ge\frac{4^k}{2\sqrt k}$: true for $k=1$ (equality), and if it holds for $k$ then $\binom{2k+2}{k+1}=\binom{2k}k\frac{2(2k+1)}{k+1}\ge\frac{4^k}{2\sqrt k}\cdot\frac{2(2k+1)}{k+1}\ge\frac{4^{k+1}}{2\sqrt{k+1}}$, the last step being $(2k+1)^2\ge4k(k+1)$. Hence each factor is $\ge\frac{16^k(2k+1)}{4k}>\frac{16^k}2$, and
\[
\begin{gathered}
\frac{c_{2n}}{c_n^4}>2^{-(n-1)}\prod_{k=1}^{n-1}16^k=2^{-(n-1)}\,4^{n(n-1)},\\
\det H_n<2^{\,n-1}4^{-n(n-1)}=2^{-(n-1)(2n-1)}\le4^{-(n-1)^2}.
\end{gathered}
\]
This decay is why the Hilbert matrix is the classic example of an invertible matrix whose numerical inversion is unreliable in ordinary floating-point arithmetic at sufficiently large orders: its determinant, the product of its eigenvalues, is astronomically small while its largest eigenvalue stays below $\pi$ (it is about $1.27$ for $n=2$ and $1.80$ for $n=12$), so some eigenvalues are tiny and dividing by them amplifies rounding errors (Figure~\ref{fig:ellipse}, right).
\end{solution}

\subsection*{Notes on the challenges}

(1) The theorem is Perron's, from 1907, extended by Frobenius in 1912 to non-negative matrices that are \emph{irreducible} (the graph of non-zero entries is strongly connected), where the eigenvector stays positive but other eigenvalues of modulus $r$ may appear (a cyclic permutation matrix has all $n$th roots of unity). A proof for positive matrices in the spirit of Exercise~\ref{ex:stochastic}: show that $r=\max\{t\ge0:\ Ax\ge tx\text{ for some }x\ge0,\ x\ne0\}$ is attained, that a maximiser is a positive eigenvector, and that any other eigenvalue has smaller modulus by applying the inequality $|Az|\le A|z|$ entrywise. Horn and Johnson, Chapter~8, is the reference. The damping factor $0.15$ of PageRank makes the Google matrix positive, so that none of the subtleties of the reducible case arise.

(2) Normalise a Hadamard matrix so that its first row is all ones (multiply columns by $\pm1$); orthogonality of rows $1$ and $2$ says row $2$ has $n/2$ entries $+1$, so $n$ is even; then counting the columns where rows $2$ and $3$ agree in sign, using orthogonality of the three pairs, gives $n\equiv0\pmod4$ for $n\ge3$. The conjecture that Hadamard matrices exist for every $n\equiv0\pmod4$ is from Paley's paper of 1933 and remains open; order $428$ was resolved in 2005. In August 2026, Alp\"oge, Voinov and Reynolds-Haertle announced the twelve previously unknown orders below $2000$, including $668$; Egan and coauthors verified them in their October 2026 survey, which lists $2396$ as the smallest still-open order.

(3) For $n=2$, writing the nilpotent $N=\begin{pmatrix}0&1\\0&0\end{pmatrix}$ as $S_1S_2$ is already a nice puzzle (try $S_1=\begin{pmatrix}0&1\\1&0\end{pmatrix}$, and $S_2$ follows); the general theorem is Frobenius's, from 1910, and the clean proof uses the rational canonical form of Volume~2: a companion matrix is visibly a product of two symmetric matrices, and the property passes to direct sums and to similar matrices (if $A=S_1S_2$ then $P^{-1}AP=(P^{-1}S_1P^{-\mathsf T})\cdot(P^{\mathsf T}S_2P)$, both symmetric). A.~J.~Bosch's note in the \emph{Monthly} gives a two-page account.

\section*{Sources and further reading}

\begin{itemize}\raggedright
\item R.~Egan, I.~S.~Kotsireas, S.~London and P.~\'O~Cath\'ain, \emph{Recent progress in the construction of Hadamard matrices}, arXiv:2610.27175, 19 October 2026. Verification of the new orders, including $668$, and the remaining frontier.
\item S.~Axler, \emph{Linear Algebra Done Right}, 4th ed., Springer, 2024 (open access). Eigenvalues, inner products and the spectral theorem, with determinants at the end.
\item P.~D.~Lax, \emph{Linear Algebra and Its Applications}, 2nd ed., Wiley, 2007. Short, deep, and full of applications: Rayleigh quotients, positive definite and positive matrices.
\item R.~A.~Horn and C.~R.~Johnson, \emph{Matrix Analysis}, 2nd ed., Cambridge University Press, 2013. The reference for everything in this chapter and much more (Chapter~4 for Rayleigh quotients, Chapter~8 for Perron--Frobenius).
\item T.~Muir, \emph{The Theory of Determinants in the Historical Order of Development}, vol.~1, Macmillan, 1906. Seki's manuscript of 1683, Leibniz's letter to l'H\^opital of 1693, Cramer's rule (1750) and Cauchy's memoir of 1812 (printed in \emph{J.\ \'Ecole Polytechnique} 10, 1815), all quoted at length.
\item A.~Cayley, \emph{A memoir on the theory of matrices}, Philos.\ Trans.\ R.\ Soc.\ London 148 (1858), 17--37. Matrices as objects, and the theorem of Theorem~\ref{thm:cayleyhamilton}, verified for $n\le3$; the general proof is G.~Frobenius, \emph{Ueber lineare Substitutionen und bilineare Formen}, J.\ reine angew.\ Math.\ 84 (1878), 1--63. The word ``matrix'' is from J.~J.~Sylvester, Philos.\ Mag.\ 37 (1850), 363--370.
\item A.-L.~Cauchy, \emph{Sur l'\'equation \`a l'aide de laquelle on d\'etermine les in\'egalit\'es s\'eculaires des mouvements des plan\`etes}, Exercices de math\'ematiques 4 (1829). Real eigenvalues of a symmetric form, before matrices. The determinant of Exercise~\ref{ex:cauchydet} is in his \emph{M\'emoire sur les fonctions altern\'ees et sur les sommes altern\'ees}, Exercices d'analyse et de physique math\'ematique 2 (1841), 151--159.
\item D.~Hilbert, \emph{Ein Beitrag zur Theorie des Legendre'schen Polynoms}, Acta Math.\ 18 (1894), 155--159. The matrix $H_n$; and M.-D.~Choi, \emph{Tricks or treats with the Hilbert matrix}, Amer.\ Math.\ Monthly 90 (1983), 301--312, for its exact inverse and why it is so badly conditioned.
\item J.~Hadamard, \emph{R\'esolution d'une question relative aux d\'eterminants}, Bull.\ Sci.\ Math.\ 17 (1893), 240--246. The inequality of Exercise~\ref{ex:hadamard}; R.~E.~A.~C.~Paley, \emph{On orthogonal matrices}, J.\ Math.\ Phys.\ 12 (1933), 311--320, for the conjecture; H.~Kharaghani and B.~Tayfeh-Rezaie, \emph{A Hadamard matrix of order 428}, J.\ Combin.\ Des.\ 13 (2005), 435--440.
\item O.~Perron, \emph{Zur Theorie der Matrices}, Math.\ Ann.\ 64 (1907), 248--263; G.~Frobenius, \emph{\"Uber Matrizen aus nicht negativen Elementen}, Sitzungsber.\ Preuss.\ Akad.\ Wiss.\ (1912), 456--477. Challenge (1).
\item A.~J.~Bosch, \emph{The factorization of a square matrix into two symmetric matrices}, Amer.\ Math.\ Monthly 93 (1986), 462--464. Challenge (3), with Frobenius's 1910 theorem.
\item S.~Brin and L.~Page, \emph{The anatomy of a large-scale hypertextual Web search engine}, Computer Networks and ISDN Systems 30 (1998), 107--117; and K.~Bryan and T.~Leise, \emph{The \$25{,}000{,}000{,}000 eigenvector: the linear algebra behind Google}, SIAM Review 48 (2006), 569--581. PageRank, as in Exercise~\ref{ex:stochastic}.
\item C.~F.~Gauss, \emph{Theoria motus corporum coelestium}, Perthes and Besser, 1809. Least squares and the orbit of Ceres; Legendre had printed the method first, in 1805, and the quarrel over priority is told in S.~M.~Stigler, \emph{The History of Statistics}, Harvard University Press, 1986, Chapter~1.
\item R.~Simson, \emph{An explication of an obscure passage in Albert Girard's Commentary upon Simon Stevin's works}, Philos.\ Trans.\ 48 (1753), 368--377. The identity $F_{n+1}F_{n-1}-F_n^2=(-1)^n$, which Cassini had stated in the \emph{Histoire de l'Acad\'emie royale des sciences} for 1680; and W.~Weaver, \emph{Lewis Carroll and a geometrical paradox}, Amer.\ Math.\ Monthly 45 (1938), 234--236, for the $64=65$ puzzle in Carroll's papers.
\end{itemize}
